\documentclass{article}
\usepackage{readmasters}
\begin{document}

\origpage{285}
\section*{Complément à l'Analysis situs;}

Par M.~H.~Poincaré, à Paris.

\emph{Adunanza del 26 marzo 1899.}

\subsection*{\S~I. Introduction.}

Dans le Journal de l'École Polytechnique (volume du centenaire de la fondation de l'École, 1894) j'ai publié un mémoire intitulé \emph{Analysis situs}, où j'étudie les variétés de l'espace à plus de trois dimensions et les propriétés des nombres de \emph{Betti}. C'est à ce mémoire que se rapporteront les renvois que je serai amené à faire fréquemment dans la suite, en mentionnant seulement le titre \emph{Analysis situs}.

Dans ce mémoire se trouve énoncé le théorème suivant~: \emph{Pour toute variété fermée, les nombres de Betti également distants des extrêmes, sont égaux}.

Le même théorème a été énoncé par M.~\emph{Picard} dans sa \emph{Théorie des fonctions algébriques de deux variables}.

M.~\emph{Heegaard} vient de revenir sur ce même problème dans un travail très remarquable, publié en langue danoise, sous le titre «~\emph{Forstudier til en topologisk teori för de algebraiske Sammenhäng}~» (Copenhague, det Nordiske Forlag Ernst Bojesen, 1898). D'après lui, le théorème en question est inexact et les démonstrations sont sans valeur.

\origpage{286}

Avant d'examiner les objections de M.~\emph{Heegaard}, il convient de faire une distinction. Il y a deux manières de définir les nombres de \emph{Betti}.

Considérons une variété $V$ que je supposerai, par exemple, fermée; soient $v_{1}, v_{2}, \ldots, v_{n}$, $n$ variétés à $p$ dimensions, faisant partie de $V$. Je suppose qu'on ne puisse pas trouver de variété à $p+1$ dimensions, faisant partie de $V$ et dont $v_{1}, v_{2}, \ldots, v_{n}$ constituent la frontière complète; mais que, si on leur adjoint une $(n+1)^{\text{ème}}$ variété à $p$ dimensions, que j'appellerai $v_{n+1}$ et qui fera partie de $V$, on puisse trouver une variété à $p+1$ dimensions, faisant partie de $V$, dont $v_{1}, v_{2}, \ldots, v_{n}, v_{n+1}$ constituent la frontière complète et \emph{cela de quelque manière que l'on ait choisi la $(n+1)^{\text{ème}}$ variété $v_{n+1}$}. Dans ce cas, on dit que le nombre de \emph{Betti} est égal à $n+1$ pour les variétés à $p$ dimensions.

C'est la définition adoptée par \emph{Betti}.

Mais on peut donner une seconde définition.

Supposons que l'on puisse trouver dans $V$ une variété à $p+1$ dimensions, dont $v_{1}, v_{2}, \ldots, v_{n}$ constituent la frontière complète; j'exprimerai ce fait par la relation suivante~:
\[ v_{1} + v_{2} + \ldots + v_{n} \backsim 0, \]
que j'appellerai une \emph{homologie}.

Il pourra se faire que sur la frontière complète de notre variété à $p+1$ dimensions, une même variété $v_{1}$ se retrouve plusieurs fois; dans ce cas, elle figurera dans le premier membre de l'homologie avec un coefficient, qui devra être un nombre entier.

D'après cette définition, on peut additionner les homologies, les soustraire les unes des autres, les multiplier par un nombre entier.

Nous \emph{conviendrons} également qu'il est permis de diviser une homologie par un nombre entier, quand tous les coefficients sont divisibles par cet entier. Par conséquent, s'il y a une variété à $p+1$ dimensions, dont la frontière complète sera constituée par 4 fois la variété $v_{1}$, nous conviendrons qu'on peut écrire non seulement l'homologie~:
\[ 4v_{1} \backsim 0, \]
\origpage{287}
mais encore l'homologie~:
\[ v_{1} \backsim 0; \]
de sorte que cette homologie signifie qu'il y a des variétés à $p+1$ dimensions, qui admettent pour frontière complète la variété $v_{1}$ ou un certain nombre de fois cette variété.

L'homologie
\[ 2v_{1} + 3v_{2} \backsim 0 \]
signifie qu'il y a des variétés à $p+1$ dimensions, qui ont pour frontière complète 2 fois $v_{1}$ et 3 fois $v_{2}$, ou 4 fois $v_{1}$ et 6 fois $v_{2}$, ou 6 fois $v_{1}$ et 9 fois $v_{2}$, etc.

Telles sont les conventions que j'ai adoptées dans l'\emph{Analysis situs}, page 19.

Je dirai que plusieurs variétés sont indépendantes, si elles ne sont pas liées par aucune homologie à coefficients entiers.

Si alors il y a $n$ variétés indépendantes à $p$ dimensions, le nombre de \emph{Betti}, d'après la seconde définition, est égal à $n+1$.

Cette seconde définition, qui est celle que j'ai adoptée dans l'\emph{Analysis situs}, ne concorde pas avec la première.

Le théorème énoncé plus haut, et critiqué par M.~\emph{Heegaard}, est vrai pour les nombres de \emph{Betti}, définis de la seconde manière, et faux pour les nombres de \emph{Betti}, définis de la première manière.

C'est ce que prouve l'exemple cité par M.~\emph{Heegaard}, page 86.

Si l'on adopte la première définition, on a~:
\[ P_{1} = 2, \qquad P_{2} = 1 \]
et par conséquent
\[ P_{2} < P_{1}. \]

Si l'on adopte, au contraire, la seconde définition, on trouve~:
\[ P_{1} = 1, \qquad P_{2} = 1, \]
et par conséquent
\[ P_{2} = P_{1}, \]
conformément au théorème énoncé.

\origpage{288}

C'est ce que prouvait également un exemple que j'ai cité, moi-même, dans l'\emph{Analysis situs}. C'est le 3\textsuperscript{ème} exemple, page 51. Mais en renvoyant à cet exemple et à cette page de l'\emph{Analysis situs}, je dois signaler une faute d'impression, qui pourrait gêner le lecteur${}^{(*)}$.

Au lieu de
\[ ABB'A' \equiv DD'C'C, \]
il faut lire
\[ ABB'A' \equiv C'CDD'; \]
au lieu de
\[ AB \equiv B'D' \equiv C'A', \]
il faut lire
\[ AB \equiv B'D' \equiv C'C. \]

Nous avons formé (page 66) les équivalences fondamentales, qui s'écrivent de la façon suivante~:
\[ 2C_{1} \equiv 2C_{2} \equiv 2C_{3} \equiv 0, \qquad 4C_{1} \equiv 0; \]
nous pouvons en déduire les homologies~:
\[ 4C_{1} \backsim 4C_{2} \backsim 4C_{3} \backsim 0. \]

Comme, d'après notre convention, on peut diviser ces homologies par 4, nous arrivons au système suivant d'homologies fondamentales~:
\[ C_{1} \backsim C_{2} \backsim C_{3} \backsim 0. \]

Si alors $P_{1}$ et $P_{2}$ sont les nombres de \emph{Betti}, \emph{définis de la seconde manière}, on trouve
\[ P_{1} = P_{2} = 1. \]

\textbf{(*)} Je profite de l'occasion pour signaler un autre erratum de l'\emph{Analysis situs}. Au bas de la page 102, il faut lire~: «~ou plus précisément les deux $v_{q+1}$, auxquelles $v_{q}$ appartiendra, appartiendront aux mêmes $v_{q+1}$\ednote{Printed $v_{q+1}$; the sequence $v_{q+1}$, $v_{q+3}$, \ldots, $v_{p}$ requires $v_{q+2}$ here, an apparent misprint.}, aux mêmes $v_{q+3}$, \ldots, aux mêmes $v_{p}$; de telle façon que la suppression de $v_{q}$ et l'annexion mutuelle des deux $v_{q+1}$, ne changera rien aux $v_{q+2}$, aux $v_{q+3}$, \ldots, aux $v_{p}$~».

\origpage{289}

Mais l'égalité entre les nombres $P_{1}$ et $P_{2}$ ne subsisterait pas si on avait adopté la première définition, qui est celle de \emph{Betti}; nous aurions toujours $P_{2} = 1$, mais nous n'aurions plus $P_{1} = 1$~:

En effet, il n'y a pas de variété à deux dimensions qui ait pour frontière complète la ligne fermée $C_{1}$, sans quoi nous aurions l'équivalence $C_{1} \equiv 0$.

Ce qui est vrai seulement, c'est qu'il y a une variété à deux dimensions, admettant pour frontière 4 fois la ligne $C_{1}$. Donc $P_{1}$ n'est pas égal à 1.

Revenons au théorème d'après lequel les nombres de \emph{Betti}, également distants des extrêmes, sont égaux.

La démonstration que j'en ai donnée dans l'\emph{Analysis situs}, semble s'appliquer également bien aux deux définitions des nombres de \emph{Betti}; elle doit donc avoir un point faible, puisque les exemples qui précèdent, montrent suffisamment que le théorème n'est pas vrai pour la première définition.

M.~\emph{Heegaard} s'en est bien rendu compte; mais je ne crois pas que sa première objection soit fondée.

Après avoir cité la façon dont je définis les variétés $V_{1}, V_{2}, \ldots, V_{p}$ (\emph{Analysis situs}, page 43), par les équations $\Phi = 0$, $F''_{i} = 0$, il ajoute (page 70) «~\emph{Enhver af Mangfoldighederne $V$ skulde altsaa kunne være den fulstaendige Skoering mellem $p$ Mangfoldigheder af $h - 1$ Dimensioner i $U$}~». Cela n'est pas exact, car, outre mes égalités, j'ai un certain nombre d'inégalités, que j'ai introduites au début du mémoire et que j'ai négligé d'écrire de nouveau dans la suite; mes variétés ne sont donc pas des intersections complètes.

La seconde objection est, au contraire, fondée. «~\emph{Naar omvendt}, dit M.~\emph{Heegaard}, \emph{Homologien $\displaystyle\sum V_{i} \backsim 0$ ikke finder Sted, saa i $U'$ kan legges en lukket Kurve $V'$, saa at}
\[ \sum N(V', V_{i}) \neq 0 \]
\emph{men det er ikke sikkert, at denne Kurve kan udskæres af nogen Mangfoldighed $V$}~». C'est là, en effet, le véritable point faible de la démonstration.

\origpage{290}

Il est donc nécessaire de revenir sur la question, et c'est l'objet du présent travail.

Souvent, pour simplifier les démonstrations, j'ai envisagé seulement le cas des variétés fermées à 3 dimensions, contenues dans l'espace à 4 dimensions. On pourrait facilement les étendre au cas général.

J'envisage, donc, dans la suite, une variété $V$ fermée, mais pour calculer ses nombres de \emph{Betti}, je la suppose divisée en variétés plus petites, de façon à former un polyèdre, au sens donné à ce mot à la page 100 de l'\emph{Analysis situs}.

\subsection*{\S~II. Schéma d'un polyèdre.}

Considérons, donc, comme à la page 100 de l'\emph{Analysis situs}, un polyèdre à $p$ dimensions, c'est-à-dire une variété $V$ à $p$ dimensions, divisée en variétés $v_{p}$; les frontières des $v_{p}$ seront les $v_{p-1}$, celles des $v_{p-1}$ seront les $v_{p-2}$, \ldots, celles des $v_{1}$ (arêtes) seront les $v_{0}$ (sommets).

J'appellerai $\alpha_{i}$ le nombre des $v_{i}$.

Soient $a^{q}_{1}, a^{q}_{2}, \ldots, a^{q}_{\alpha_{q}}$ les différentes $v_{q}$.

Soit $a^{q}_{1}$ une des variétés $v_{q}$ et $a^{q-1}_{1}$ une des variétés $v_{q-1}$, qui lui sert de frontière. Étudions les rapports de $a^{q}_{1}$ et de $a^{q-1}_{1}$.

Soient
\[ F_{1} = F_{2} = \ldots = F_{n-q} = F_{n-q+1} = 0, \qquad \varphi_{\gamma} > 0 \tag{1} \]
les égalités et les inégalités qui définissent $a^{q-1}_{1}$, d'après la première définition des variétés (\emph{Analysis situs}, page 4).

Les relations qui définissent $a^{q}_{1}$ pourront se mettre sous la forme~:
\[ F_{1} = F_{2} = \ldots = F_{n-q} = 0, \qquad F_{n-q+1} > 0, \qquad \varphi_{\gamma} > 0. \tag{2} \]

Dans ce cas, nous dirons que la relation de $a^{q}_{1}$ et de $a^{q-1}_{1}$ est \emph{directe}.

Cette relation deviendrait \emph{inverse}, si l'une de ces deux variétés était remplacée par la variété opposée; elle redeviendrait directe, si chacune des deux variétés était remplacée par la variété opposée.

\origpage{291}

On sait qu'une variété est remplacée par la \emph{variété opposée} (\emph{Analysis situs}, page 15), quand on permute deux des fonctions $F$ (qui, égalées à zéro, donnent les équations qui définissent la variété), ou qu'on change le signe de l'une d'elles.

Ainsi les deux variétés
\[ \begin{aligned} &F_{1} = F_{2} = F_{3} = 0; \qquad F_{1} = F_{2} = 0, \qquad F_{3} > 0;\\ &F_{1} = F_{2} = F_{3} = 0; \qquad F_{1} = F_{3} = 0, \qquad F_{2} < 0;\\ &F_{1} = F_{2} = F_{3} = 0; \qquad F_{2} = F_{3} = 0, \qquad F_{1} > 0; \end{aligned} \]
sont en relation directe; tandis que les deux variétés
\[ \begin{aligned} &F_{1} = F_{2} = F_{3} = 0; \qquad F_{1} = F_{2} = 0, \qquad F_{3} < 0;\\ &F_{1} = F_{2} = F_{3} = 0; \qquad F_{1} = F_{3} = 0, \qquad F_{2} > 0; \end{aligned} \]
sont en relation inverse.

Cela posé, soit $\varepsilon^{q}_{i,j}$ un nombre qui sera égal à 0, si $a^{q-1}_{j}$ n'est pas frontière de $a^{q}_{i}$; à $+1$, si $a^{q-1}_{j}$ est frontière de $a^{q}_{i}$ et en relation directe avec $a^{q}_{i}$; et, enfin, à $-1$, si $a^{q-1}_{j}$ est frontière de $a^{q}_{i}$, mais en relation inverse avec $a^{q}_{i}$.

Nous conviendrons d'écrire la \emph{congruence}
\[ a^{q}_{i} \equiv \sum_{j} \varepsilon^{q}_{i,j} a^{q-1}_{j}, \tag{3} \]
qui nous fait connaitre les frontières de $a^{q}_{i}$.

L'ensemble des congruences (3), relatives aux différentes $v_{p}, v_{p-1}, \ldots, v_{0}$ de $V$, constitue ce qu'on peut appeler le \emph{schéma d'un polyèdre}.

On peut se poser deux questions~:

1° Un schéma étant donné, existera-t-il toujours un polyèdre, qui y corresponde?

2° Deux polyèdres qui ont même schéma, sont-ils homéomorphes?

Sans aborder, pour le moment, ces deux questions, cherchons
\origpage{292}
quelques-unes des conditions auxquelles doit satisfaire un schéma, pour qu'un polyèdre y puisse correspondre.

Considérons l'une des $v_{p-1}$, $a^{p-1}_{i}$ par exemple; cette variété devra séparer, l'une de l'autre, deux des $v_{p}$ et deux seulement; de sorte que, parmi les nombres
\[ \varepsilon^{p}_{i,1}, \]\ednote{The indices do not agree as printed: the text names the face $a^{p-1}_{i}$, while this display and the sum $\displaystyle\sum_{i} \varepsilon^{p}_{i,1}$ on p.~293 treat the face as $a^{p-1}_{1}$. Either the text's index or the display's is a misprint.}
il y en aura un qui sera égal à $+1$, un qui sera égal à $-1$ et tous les autres seront égaux à 0.

Ce n'est pas tout; envisageons l'une quelconque des $v_{q}$, $a^{q}_{i}$ par exemple, et une quelconque des $v_{q-2}$, $a^{q-2}_{k}$ par exemple.

De deux choses, l'une~: ou bien $a^{q-2}_{k}$ n'appartiendra pas à $a^{q}_{i}$, et dans ce cas tout\ednote{Printed ``tout''; read ``tous''.} les produits
\[ \varepsilon^{q}_{i,j} \varepsilon^{q-1}_{j,k} \tag{4} \]
seront nuls, car si $a^{q-1}_{j}$ n'appartient pas à $a^{q}_{i}$, le premier facteur est nul; si, au contraire, $a^{q-1}_{j}$ appartient à $a^{q}_{i}$, la variété $a^{q-2}_{k}$ ne peut pas appartenir à $a^{q-1}_{j}$ (sans quoi, elle appartiendrait à $a^{q}_{i}$, contrairement à l'hypothèse), et le second facteur doit être nul.

Ou bien $a^{q-2}_{k}$ appartiendra à $a^{q}_{i}$; mais alors nous pourrons raisonner sur la variété $a^{q}_{i}$, comme nous raisonnions tout à l'heure sur la variété $V$, et nous conclurons que $a^{q-2}_{k}$ doit séparer, l'une de l'autre, deux des variétés $v_{q-1}$, qui appartiennent à $a^{q}_{i}$, et deux seulement~: soient $a^{q-1}_{1}$ et $a^{q-1}_{2}$.

Parmi les produits (4) il n'y en aura que deux qui ne seront pas nuls, à savoir
\[ \varepsilon^{q}_{i,1} \varepsilon^{q-1}_{1,k}, \qquad \varepsilon^{q}_{i,2} \varepsilon^{q-1}_{2,k}. \]

Pour tous les autres, en effet, ou bien $a^{q-1}_{j}$ n'appartiendra pas à $a^{q}_{i}$, ou bien $a^{q-2}_{k}$ n'appartiendra pas à $a^{q-1}_{j}$.

Ces deux produits sont d'ailleurs égaux~: l'un à $+1$, l'autre à $-1$.

On aura donc dans tous les cas~:
\[ \sum_{j} \varepsilon^{q}_{i,j} \varepsilon^{q-1}_{j,k} = 0. \tag{5} \]

\origpage{293}

Nous avons de même
\[ \sum_{i} \varepsilon^{p}_{i,1} = 0, \]
et, plus généralement, quel que soit $k$~:
\[ \sum_{i} \varepsilon^{p}_{i,k} = 0. \tag{5 \textit{bis}} \]

La relation (5~\emph{bis}) peut être regardée, à un certain point de vue, comme un cas particulier de la relation (5).

Soit $P$ la portion de l'espace à $p+1$ dimensions, limitée par le polyèdre $V$; alors la frontière complète de $P$ se composera des diverses variétés $v_{p}$, qui, par leur ensemble, forment $V$; nous pourrons donc écrire, au sens de la congruence (3),
\[ P \equiv \sum_{i} a^{p}_{i}, \tag{3 \textit{bis}} \]
ou encore
\[ P \equiv \sum_{i} \varepsilon^{p+1}_{0,i} a^{p}_{i}, \]
où les nombres $\varepsilon^{p+1}_{0,i}$ seront tous, par définition, égaux à 1.

A ce compte, la relation (5~\emph{bis}), qui peut s'écrire
\[ \sum_{i} \varepsilon^{p+1}_{0,i} \varepsilon^{p}_{i,k} = 0, \]
n'est plus qu'un cas particulier de la relation (5).

Nous avons, ensuite, chaque $v_{1}$ qui a pour limites deux $v_{0}$, et deux seulement, ce qui nous donne des congruences (3) de la forme~:
\[ a^{1}_{i} \equiv a^{0}_{j} - a^{0}_{k}, \]
et une relation analogue à (5) et (5~\emph{bis})~:
\[ \sum_{j} \varepsilon^{1}_{i,j} = 0, \]
qui rentrerait encore dans la forme (5), en convenant de faire tous les $\varepsilon^{0}$ égaux à $+1$.

\origpage{294}

D'autre part, envisageons l'une des $a^{q}_{i}$, toutes les $a^{q+1}_{j}$ auxquelles elle sert de frontière; toutes les $a^{q+2}_{k}$, auxquelles ces $a^{q+1}_{j}$ servent de frontières et ainsi de suite. L'ensemble de toutes ces variétés constituera ce que nous avons appelé un \emph{aster} (\emph{Analysis situs}, page 106).

Nous avons vu (loc.\ cit., page 107) que le polyèdre qui correspond à un aster, doit être simplement connexe. Ainsi une condition pour qu'un polyèdre puisse correspondre à un schéma donné, c'est que les polyèdres qui correspondent aux différents asters, d'après la convention de la page 107 de l'\emph{Analysis situs}, soient tous simplement connexes.

Considérons maintenant une des $a^{q}_{i}$, toutes les $a^{q-1}_{j}$ qui lui servent de frontières, les $a^{q-2}_{k}$ qui servent de frontières à ces $a^{q-1}_{j}$ et ainsi de suite. Cet ensemble de variétés constituera un polyèdre à $q$ dimensions; nous supposerons que ce polyèdre soit simplement connexe.

Ce n'est plus là une condition nécessaire pour qu'un polyèdre puisse correspondre au schéma; c'est simplement une condition que, sauf avis contraire, nous supposerons remplie.

Pour éclairer ces définitions par quelques exemples, voyons d'abord quel est le schéma du tétraèdre généralisé, défini à la page 105 (\emph{Analysis situs}).

Les faces de ce tétraèdre seront définies par les $n+1$ équations~:
\[ \begin{aligned} &x_{1} = 0, \qquad x_{2} = 0, \qquad \ldots, \qquad x_{n} = 0\\ &x_{1} + x_{2} + \ldots + x_{n} = 1. \end{aligned} \tag{6} \]

On obtiendra les $a^{q}_{i}$ en supprimant $q+1$ de ces équations; pour définir le sens de la variété $a^{q}_{i}$, nous supposerons qu'on supprime ces $q+1$ équations, mais sans changer l'ordre des $n-q$ équations restantes.

Cela posé, considérons la relation de $a^{q}_{i}$ et de $a^{q-1}_{j}$ et cherchons à déterminer le nombre $\varepsilon^{q}_{i,j}$.

D'abord, pour que $a^{q-1}_{j}$ appartienne à $a^{q}_{i}$, il faut que $a^{q-1}_{j}$ soit définie par les $n-q$ équations qui définissent $a^{q}_{i}$, auxquelles on devra adjoindre une $(n-q+1)^{\text{ème}}$ équation, prise parmi les équations (6). S'il n'est pas ainsi, le nombre $\varepsilon^{q}_{i,j}$ sera nul.

\origpage{295}

Supposons, donc, que $a^{q-1}_{j}$ soit obtenu en supprimant les $q$ équations qui occupent les
\[ \alpha_{1}, \qquad \alpha_{2}, \qquad \ldots, \qquad \alpha_{q} \]
rangs.

Supposons que $a^{q}_{i}$ soit obtenue en supprimant, \emph{en outre}, la $\beta^{\text{ème}}$ équation; alors le nombre $\varepsilon^{q}_{i,j}$, dont la valeur absolue sera toujours égale à 1, aura même signe que le produit~:
\[ (\beta - \alpha_{1})(\beta - \alpha_{2}) \ldots (\beta - \alpha_{q}). \]

Il est aisé de vérifier alors que la relation (5) a lieu.

Considérons, en effet, la variété $a^{q-2}_{k}$, obtenue en supprimant les équations de rang $\alpha_{1}, \alpha_{2}, \ldots, \alpha_{q-1}$ et la variété $a^{q}_{i}$, obtenue en supprimant, en outre, les équations de rang $\beta$ et $\gamma$. (Il est clair que si $a^{q}_{i}$ ne s'obtenait pas en supprimant les mêmes équations que pour $a^{q-2}_{k}$, plus deux autres, tous les produits $\varepsilon^{q}_{i,j} \varepsilon^{q-1}_{j,k}$ seraient nuls).

Dans ce cas, tous ces produits seront encore nuls, sauf deux~:
\[ \varepsilon^{q}_{i,1} \varepsilon^{q-1}_{1,k} \qquad \text{et} \qquad \varepsilon^{q}_{i,2} \varepsilon^{q-1}_{2,k}, \]
qui correspondront aux deux variétés $a^{q-1}_{1}$ et $a^{q-2}_{2}$\ednote{Printed $a^{q-2}_{2}$; the context requires $a^{q-1}_{2}$, an apparent misprint.}, obtenues respectivement en supprimant les équations de rang $\alpha_{1}, \alpha_{2}, \ldots, \alpha_{q-1}, \beta$ et celles de rang $\alpha_{1}, \alpha_{2}, \ldots, \alpha_{q-1}, \gamma$.

Alors les quatre nombres
\[ \varepsilon^{q}_{i,1}, \qquad \varepsilon^{q-1}_{1,k}, \qquad \varepsilon^{q}_{i,2}, \qquad \varepsilon^{q-1}_{2,k} \]
auront respectivement même signe que~:
\[ \begin{aligned} (\gamma - \beta)(\gamma - \alpha_{1})(\gamma - \alpha_{2}) \ldots (\gamma - \alpha_{q-1})&,\\ (\beta - \alpha_{1})(\beta - \alpha_{2}) \ldots (\beta - \alpha_{q-1})&,\\ (\beta - \gamma)(\beta - \alpha_{1})(\beta - \alpha_{2}) \ldots (\beta - \alpha_{q-1})&,\\ (\gamma - \alpha_{1})(\gamma - \alpha_{2}) \ldots (\gamma - \alpha_{q-1})&. \end{aligned} \]

\origpage{296}

On vérifie ainsi que les deux produits, qui ne sont pas nuls, sont égaux et de signe contraire. C.~Q.~F.~D.${}^{(*)}$

\subsection*{\S~III. Nombres de Betti réduits.}

Je m'en vais chercher, maintenant, les nombres de \emph{Betti}, relatifs à un polyèdre, mais afin d'éviter l'équivoque, dont j'ai signalé plus haut la possibilité, je conviendrai de définir ces nombres \emph{de la seconde manière}, c'est-à-dire que $P_{q} - 1$ sera le nombre de variétés fermées à $q$ dimensions, que l'on peut tracer sur notre polyèdre $V$ et qui sont \emph{linéairement indépendantes}, je veux dire, qui ne sont liées par aucune homologie à coefficients entiers, au sens de l'\emph{Analysis situs}, page 18.

Mais je me proposerai d'abord de déterminer le nombre $P'_{q} - 1$ des variétés à $q$ dimensions, fermées et linéairement indépendantes, que l'on peut tracer sur notre polyèdre $V$, mais \emph{en nous bornant à celles qui sont, des combinaisons des variétés $v_{q}$}.

Le nombre $P'_{q}$ sera, alors, ce que j'appellerai le \emph{nombre de Betti réduit}.

Les variétés à $q$ dimensions, qui sont des combinaisons des $v_{q}$, pourront évidemment être représentées par~:
\[ \sum_{i} \lambda_{i} a^{q}_{i}, \]
les $\lambda_{i}$ étant des coefficients entiers et les lettres $a^{q}_{i}$ continuant à désigner les différentes variétés $v_{q}$.

Quelle est d'abord la condition pour que la variété $\displaystyle\sum_{i} \lambda_{i} a^{q}_{i}$ soit fermée?

Pour cela, cherchons quelles sont les variétés $v_{q-1}$, qui forment les frontières de cette variété. Pour les trouver, il suffit évidemment de remplacer $a^{q}_{i}$ par sa valeur donnée par la congruence (3).

\textbf{(*)} Le polyèdre ainsi défini, ainsi que tout polyèdre à $n$ dimensions, limité par $n+1$ variétés planes, s'appellera \emph{tétraèdre généralisé rectiligne}. J'appellerai \emph{tétraèdre généralisé} toute variété homéomorphe à un tétraèdre généralisé rectiligne.

\origpage{297}

Cet ensemble de variétés frontières sera donc donné par la formule~:
\[ \sum_{i} \sum_{j} \lambda_{i} \varepsilon^{q}_{i,j} a^{q-1}_{j}. \]

Pour que la variété $\displaystyle\sum_{i} \lambda_{i} a^{q}_{i}$ soit fermée, il suffit, donc, que l'on ait identiquement~:
\[ \sum_{i} \sum_{j} \lambda_{i} \varepsilon^{q}_{i,j} a^{q-1}_{j} = 0, \]
c'est-à-dire que, quel que soit $j$, on ait~:
\[ \sum_{i} \lambda_{i} \varepsilon^{q}_{i,j} = 0. \]

En d'autres termes, la variété $\displaystyle\sum_{i} \lambda_{i} a^{q}_{i}$ sera fermée, si l'on a~:
\[ \sum_{i} \lambda_{i} a^{q}_{i} \equiv 0, \tag{7, $q$} \]
en vertu des congruences (3, $q$); j'appelle ainsi celles des congruences (3), qui lient les $a^{q}_{i}$ aux $a^{q-1}_{j}$.

Cherchons maintenant les homologies qui peuvent exister entre les variétés $a^{q}_{i}$. On obtiendra toutes ces homologies, en combinant celles que l'on peut obtenir de la façon suivante.

Considérons la congruence~:
\[ a^{q+1}_{k} \equiv \sum_{i} \varepsilon^{q+1}_{k,i} a^{q}_{i}, \tag{8} \]
qui, d'après la convention que nous venons de faire,.\ednote{A stray point follows the comma in the print.} est une congruence (3, $q + 1$); remplaçons le signe $\equiv$ par $\backsim$, et le premier membre par zéro; il viendra~:
\[ \sum_{i} \varepsilon^{q+1}_{k,i} a^{q}_{i} \backsim 0. \tag{9, $q$} \]

Cette homologie aura évidemment lieu, puisque, par définition, elle exprime, comme la congruence (8), que les $a^{q}_{i}$ forment la frontière complète de $a^{q+1}_{k}$.

\origpage{298}

Nous démontrerons plus loin (\S~6), qu'il n'y en a pas d'autres.

Je désigne cette homologie par (9, $q$) pour marquer qu'elle a lieu entre les $a^{q}_{i}$.

Je dis que si l'homologie (9, $q$) a lieu, la congruence
\[ \sum_{i} \varepsilon^{q+1}_{k,i} a^{q}_{i} \equiv 0 \tag{10, $q$} \]
sera une conséquence des congruences (3, $q$).

Remplaçons, en effet, les $a^{q}_{i}$ par leurs valeurs, données par ces congruences (3, $q$); il viendra~:
\[ \sum_{i} \varepsilon^{q+1}_{k,i} a^{q}_{i} \equiv \sum_{i} \sum_{j} \varepsilon^{q+1}_{k,i} \varepsilon^{q}_{i,j} a^{q-1}_{j}. \]
Le second membre est identiquement nul en vertu des relations (5).

Cela posé, soit $\alpha_{q}$ le nombre des variétés $a^{q}_{i}$; soit $\alpha'_{q}$ le nombre de ces variétés qui restent distinctes, si l'on ne regarde pas comme distinctes des variétés liées par une homologie de la forme (9, $q$); soit $\alpha''_{q}$ le nombre de ces variétés qui restent distinctes, si l'on ne regarde pas comme distinctes des variétés liées par une congruence de la forme (7, $q$).

Il résulte de ces définitions~:

1° qu'il y a $\alpha_{q} - \alpha'_{q}$ homologies distinctes de la forme (9, $q$);

2° qu'il y a $\alpha_{q} - \alpha''_{q}$ congruences distinctes de la forme (7, $q$);

3° que
\[ \alpha'_{q} \geqq \alpha''_{q}; \]
car si plusieurs $a^{q}_{i}$ sont liées par une homologie de la forme (9, $q$), elles seront liées également par la congruence (10, $q$) correspondante.

Enfin le nombre cherché $P'_{q} - 1$ est égal à
\[ \alpha'_{q} - \alpha''_{q}, \]
car les variétés fermées de la forme $\displaystyle\sum_{i} \lambda_{i} a^{q}_{i}$, réellement distinctes, sont en nombre égal à celui des congruences (7, $q$), c'est-à-dire au nombre de $\alpha_{q} - \alpha''_{q}$.

Le nombre $P'_{q} - 1$ est le nombre de ces variétés qui restent
\origpage{299}
distinctes, en ne regardant pas comme distinctes celles qui sont liées par une homologie (9, $q$). Or le nombre de ces homologies est $\alpha_{q} - \alpha'_{q}$; nous avons donc~:
\[ P'_{q} - 1 = (\alpha_{q} - \alpha''_{q}) - (\alpha_{q} - \alpha'_{q}) = \alpha'_{q} - \alpha''_{q} \qquad \text{C.~Q.~F.~D.} \]
Soient $a^{q}_{1}, a^{q}_{2}, \ldots, a^{q}_{i}$ des variétés $v_{q}$, au nombre de $i$ et
\[ \begin{aligned} a^{q}_{1} &\equiv \sum \varepsilon^{q}_{1,j} a^{q-1}_{j}, \\ a^{q}_{2} &\equiv \sum \varepsilon^{q}_{2,j} a^{q-1}_{j}, \\ &\ldots\ldots\ldots\ldots \\ a^{q}_{i} &\equiv \sum \varepsilon^{q}_{i,j} a^{q-1}_{j} \end{aligned} \]
les congruences (3) correspondantes. Formons les homologies correspondantes~:
\[ \sum \varepsilon^{q}_{1,j} a^{q-1}_{j} \backsim 0, \qquad \sum \varepsilon^{q}_{2,j} a^{q-1}_{j} \backsim 0, \qquad \ldots, \qquad \sum \varepsilon^{q}_{i,j} a^{q-1}_{j} \backsim 0. \]
La condition nécessaire et suffisante pour que ces homologies soient distinctes, c'est que l'on n'ait entre les $i$ variétés $a^{q}_{1}, a^{q}_{2}, \ldots, a^{q}_{i}$ aucune congruence de la forme~:
\[ \lambda_{1} a^{q}_{1} + \lambda_{2} a^{q}_{2} + \ldots + \lambda_{i} a^{q}_{i} \equiv 0. \]

Le nombre des homologies distinctes est donc égal au nombre des $a^{q}_{i}$ distinctes, en tenant compte des congruences (7, $q$). Donc~:
\[ \alpha_{q-1} - \alpha'_{q-1} = \alpha''_{q}, \]
ou
\[ \alpha_{q-1} = \alpha'_{q-1} + \alpha''_{q}. \]
Nous aurons d'autre part
\[ \alpha'_{0} = 1. \]
Si, en effet, on peut aller d'un sommet quelconque $a^{0}_{1}$ à un autre
\origpage{300}
sommet quelconque $a^{0}_{i}$, en suivant des arêtes (c'est-à-dire si le polyèdre est d'un seul tenant), on aura l'homologie~:
\[ a^{0}_{1} \backsim a^{0}_{i}, \]
c'est-à-dire qu'il n'y aura qu'un seul sommet distinct, en tenant compte des homologies.

Envisageons maintenant la congruence (3~\emph{bis})
\[ P \equiv \sum a^{p}_{i}; \]
l'homologie correspondante s'écrit~:
\[ \sum a^{p}_{i} \backsim 0, \]
et il n'y a pas d'autre homologie (9, $p$). Donc
\[ \alpha_{p} = \alpha'_{p} + 1. \]
De plus, le polyèdre étant d'un seul tenant, une seule des combinaisons $\displaystyle\sum \lambda_{i} a^{p}_{i}$ pourra être fermée, c'est le polyèdre lui-même dans son entier, représenté par la formule $\displaystyle\sum a^{p}_{i}$.

Nous aurons donc une seule congruence de la forme (7, $p$)~:
\[ \sum a^{p}_{i} \equiv 0. \]
Donc
\[ \alpha_{p} = \alpha''_{p} + 1, \qquad \alpha'_{p} = \alpha''_{p}. \]
Nous avons donc la série d'équations~:
\[ \begin{aligned} &\alpha'_{0} = 0, \\ &\alpha_{0} = \alpha'_{0} + \alpha''_{1}, &&\qquad \alpha'_{1} - \alpha''_{1} = P'_{1} - 1, \\ &\alpha_{1} = \alpha'_{1} + \alpha''_{2}, &&\qquad \alpha'_{2} - \alpha''_{2} = P'_{2} - 1, \\ &\ldots\ldots\ldots &&\qquad \ldots\ldots\ldots\ldots \\ &\alpha_{p-1} = \alpha'_{p-1} + \alpha''_{p}, &&\qquad \alpha'_{p-1} - \alpha''_{p-1} = P'_{p-1} - 1, \\ &\alpha_{p} = \alpha'_{p} + 1, &&\qquad \alpha'_{p} - \alpha''_{p} = 0, \end{aligned} \]
\ednote{The first row is printed $\alpha'_{0} = 0$; page 299 gives $\alpha'_{0} = 1$, which the final $\mp 1$ of the formula on page 301 requires; an apparent misprint.}
\origpage{301}
d'où l'on tire aisément~:
\[ \begin{aligned} &\alpha_{p} - \alpha_{p-1} + \alpha_{p-2} - \ldots \pm \alpha_{1} \mp \alpha_{0} = 1 - (P'_{p-1} - 1) + \ldots \\ &\qquad \ldots \mp (P'_{2} - 1) \pm (P'_{1} - 1) \mp 1, \end{aligned} \]
tout à fait analogue à la formule~:
\[ \begin{aligned} &\alpha_{p} - \alpha_{p-1} + \alpha_{p-2} - \ldots \pm \alpha_{1} \mp \alpha_{0} = 1 - (P_{p-1} - 1) + \ldots \\ &\qquad \ldots \mp (P_{2} - 1) \pm (P_{1} - 1) \mp 1, \end{aligned} \]
que nous avons trouvée dans l'\emph{Analysis situs}, page 121.

\subsection*{\S~IV. Subdivision des polyèdres.}

Considérons un polyèdre $V$, à $p$ dimensions, avec ses diverses variétés~:
\[ a^{p}_{i}, \qquad a^{p-1}_{i}, \qquad \ldots, \qquad a^{1}_{i}, \qquad a^{0}_{i}. \]

Supposons que l'on subdivise chacune des variétés $a^{p}_{i}$ en plusieurs autres, que j'appellerai les $b^{p}_{i}$; soient ensuite $b^{p-1}_{i}$ les variétés à $p - 1$ dimensions, qui servent de frontière aux $b^{p}_{i}$; soient $b^{p-2}_{i}$ les variétés à $p - 2$ dimensions, qui servent de frontières aux $b^{p-1}_{i}$; \ldots\ et, enfin, $b^{0}_{i}$ les variétés à 0 dimensions (sommets), qui servent de frontières aux $b^{1}_{i}$ (arêtes).

On aura ainsi un nouveau polyèdre $V'$, qui sera dérivé du polyèdre $V$, au sens que j'ai attaché à ce mot à la page 101 de l'\emph{Analysis situs}.

On peut supposer, d'ailleurs, que si une variété $v_{q-1}$, simplement ou multiplement connexe, sert de frontière à deux variétés $b^{q}_{j}$ et $b^{q}_{k}$, elle ne forme pas forcément une seule des variétés $b^{q-1}_{i}$, mais peut être elle-même subdivisée en plusieurs variétés $b^{q-1}_{i}$. Dans ce cas, pour reprendre la terminologie des pages 102 et 103 de l'\emph{Analysis situs}, ces variétés $b^{q-1}_{i}$ seront \emph{irrégulières} et les variétés $b^{q-2}_{i}$, qui les séparent les unes des autres, seront \emph{singulières}.

\origpage{302}

Cela posé, recherchons une classification des variétés $b^{q}_{i}$.

Si une variété $b^{q}_{i}$ ne fait pas partie d'une des variétés $a^{q}_{j}$, elle fera partie d'une des variétés $a^{q+1}_{j}$, ou d'une des variétés $a^{q+2}_{j}$, \ldots, ou, tout au moins, d'une des variétés $a^{p}_{j}$.

Peut-elle faire partie à la fois de deux variétés $a^{m}_{j}$ et $a^{m}_{k}$?

D'après la façon dont la subdivision a été supposée faite, en ajoutant toujours de nouvelles frontières, sans en supprimer jamais, cela ne pourra arriver que si ces deux variétés $a^{m}_{j}$ et $a^{m}_{k}$ sont contigües et ont une frontière commune $a^{m-1}_{h}$, et si $b^{q}_{i}$ fait partie de cette frontière $a^{m-1}_{h}$.

Je suppose alors que $b^{q}_{i}$ fasse partie de $a^{h}_{j}$, et ne fasse partie d'aucune variété $a^{m}_{k}$, où $m < h$. La variété $a^{h}_{j}$ existe toujours et l'on a $h \geqq q$; de plus, la variété $a^{h}_{j}$ est unique, c'est-à-dire que $b^{q}_{i}$ ne peut faire partie à la fois de deux variétés différentes $a^{h}_{j}$ et $a^{h}_{k}$.

Si donc je conviens de ranger dans une même classe toutes les variétés $b^{q}_{i}$, qui font partie de la variété $a^{h}_{j}$, sans faire partie d'aucune des variétés $a^{m}_{k}$, où $m < h$, toute variété $b^{q}_{i}$ fera partie d'une classe et d'une seule.

Je pourrai alors représenter $b^{q}_{i}$ par une notation à quatre indices
\[ b^{q}_{i} = B(q, h, j, k); \]
l'indice $q$ indique le nombre des dimensions de $b^{q}_{i}$; les indices $h$ et $j$ indiquent que $b^{q}_{i}$ fait partie de la classe $a^{h}_{j}$; et l'indice $k$ sert à distinguer les unes des autres les différentes variétés d'une même classe. On a $h \geqq q$.

Nous aurons, alors, pour définir le polyèdre $V$ et sa subdivision~:

1° Les congruences (3, $q$), relatives au polyèdre $V$, que j'écrirai~:
\[ a^{q}_{i} \equiv \sum_{j} \varepsilon^{q}_{i,j} a^{q-1}_{j}. \tag{3, $q$, $i$} \]

2° Les équations qui donnent la subdivision de la variété $a^{q}_{i}$~:
\[ a^{q}_{i} = \sum_{k} B(q, q, i, k). \tag{1, $q$, $i$} \]

3° Les congruences analogues aux congruences (3), mais relatives
\origpage{303}
au polyèdre $V'$; je les écrirai~:
\[ B(q, h, j, k) \equiv \sum \zeta B(q - 1, h', j', k'). \tag{2, $q$, $h$, $j$, $k$} \]

Les $\zeta$ sont des nombres égaux à $\pm 1$, ou à 0; ils dépendent des sept indices $q$, $h$, $j$, $k$, $h'$, $j'$, $k'$, de sorte que je les écrirai, quand cela sera nécessaire, sous la forme~:
\[ \zeta(q, h, j, k, h', j', k'). \]

Sous le signe $\displaystyle\sum$, les indices $h'$, $j'$, $k'$ peuvent prendre toutes les valeurs. Observons, cependant, que les $B(q - 1)$, qui servent de frontière à $B(q, h, j, k)$, doivent, comme $B(q, h, j, k)$, faire partie de $a^{h}_{j}$; mais pourront faire partie d'autres variétés $a^{h'}_{j'}$, d'un nombre moindre de dimensions, mais faisant partie de $a^{h}_{j}$. On aura donc
\[ h' \leqq h, \qquad h' \geqq q - 1. \]

D'ailleurs si $h' = h$, on aura $j' = j$.

Pour que les relations (1), (2), (3) puissent définir une véritable subdivision, elles doivent satisfaire à certaines conditions.

Les relations (1, $q$, $i$), (3, $q$, $i$) donnent
\[ \sum_{k} B(q, q, i, k) \equiv \sum \varepsilon^{q}_{i,j} a^{q-1}_{j}. \tag{$\alpha$} \]

Si dans le premier membre je remplace $B(q, q, i, k)$ par sa valeur, tirée de (2, $q$, $q$, $i$, $k$), et $a^{q-1}_{j}$ par sa valeur tirée de (1, $q - 1$, $j$), les deux membres devront devenir identique\ednote{Printed identique; the agreement with membres requires identiques.}; voilà une première condition, qui est évidente; elle n'est d'ailleurs pas suffisante.

\subsection*{\S~V. Influence de la subdivision sur les nombres de Betti réduits.}

Soit $\displaystyle\sum \alpha B(q, h, j, k)$ une combinaison des variétés $b^{q}_{i}$, qui représente une variété fermée à $q$ dimensions, de telle sorte que l'on
\origpage{304}
ait, avec nos notations,
\[ \sum \alpha B(q, h, j, k) \equiv 0, \qquad (h \geqq q) \tag{1} \]

Parmi les variétés $b^{q}_{i}$ qui figurent dans le premier membre de (1), réunissons celles qui appartiennent à une même classe. Soit
\[ S \alpha B(q, h, j, k) \]
l'ensemble de celles qui appartiennent à la classe $a^{h}_{j}$; le signe de sommation $S$ signifie, donc, qu'on ne prend que les variétés d'une même classe, tandis que le signe $\displaystyle\sum$ signifie qu'on les prend toutes.

On aura alors~:
\[ S \alpha B(q, h, j, k) \equiv \sum \beta B(q - 1, h', j', k'), \tag{2} \]
c'est-à-dire que la variété à $q - 1$ dimensions $\displaystyle\sum \beta B(q - 1, h', j', k')$ forme la frontière complète de la variété à $q$ dimensions
\[ S \alpha B(q, h, j, k). \]

Les variétés $a^{h'}_{j'}$ doivent appartenir à la frontière de $a^{h}_{j}$, ou se confondre avec $a^{h}_{j}$; en effet, $B(q - 1, h', j', k')$ appartient à $a^{h'}_{j'}$ et, d'autre part, à l'une des $B(q, h, j, k)$, qui fait lui-même partie de $a^{h}_{j}$, si donc $a^{h'}_{j'}$ ne faisait pas partie de $a^{h}_{j}$, $B(q - 1, h', j', k')$ ferait partie d'une variété $a^{m}_{k}$, partie commune à $a^{h}_{j}$ et $a^{h'}_{j'}$, et qui aurait moins de $h'$ dimensions. Cela est contraire à la définition que nous avons donnée des classes.

D'autre part $a^{h'}_{j'}$ ne peut pas se confondre avec $a^{h}_{j}$.

Soit, en effet,
\[ S_{1} \alpha_{1} B(q, h_{1}, j_{1}, k_{1}) = \sum \alpha B(q, h, j, k) - S \alpha B(q, h, j, k) \]
l'ensemble des variétés qui figurent dans le premier membre de (1), et qui n'appartiennent pas à la classe $a^{h}_{j}$; on aura évidemment~:
\[ S_{1} \alpha_{1} B(q, h_{1}, j_{1}, k_{1}) \equiv - \sum \beta B(q - 1, h', j', k'). \]

\origpage{305}

Donc $B(q - 1, h', j', k')$ doit faire partie à la fois de $a^{h'}_{j'}$, et de l'une des $B(q, h_{1}, j_{1}, k_{1})$, et, par conséquent, de l'une des $a^{h_{1}}_{j_{1}}$, \emph{différentes de $a^{h}_{j}$}. Si donc $a^{h'}_{j'}$ se confondait avec $a^{h}_{j}$,
\[ B(q - 1, h', j', k') = B(q - 1, h, j, k) \]
devrait appartenir à une variété $a^{m}_{k}$, partie commune à $a^{h}_{j}$ et $a^{h_{1}}_{j_{1}}$. De deux choses, l'une~: ou bien $a^{h}_{j}$ ne ferait pas partie de $a^{h_{1}}_{j_{1}}$, et alors on aurait encore $m < h$, ce qui serait encore contraire à la définition des classes; ou bien $a^{h}_{j}$ ferait partie de $a^{h_{1}}_{j_{1}}$, et alors on aurait $h_{1} > h$. Supposons que j'ai choisi la classe $a^{h}_{j}$, qui correspond au plus grand nombre $h$. Alors on ne pourra pas avoir $h_{1} > h$, et $a^{h'}_{j'}$ devra appartenir à la frontière de $a^{h}_{j}$.

La congruence (2) entraine l'homologie
\[ \sum \beta B(q - 1, h', j', k') \backsim 0; \tag{3} \]
comme, d'autre part, $a^{h}_{j}$ \emph{est simplement connexe} et que toutes les variétés $B(q - 1, h', j', k')$ sont situées sur la frontière de $a^{h}_{j}$, le premier membre de (3), représentant une variété fermée à $q - 1$ dimensions, située sur cette frontière, formera la frontière complète d'une variété à $q$ dimensions
\[ \sum \gamma B(q, h'', j'', k''), \]
également située sur la frontière de $a^{h}_{j}$. (Il y aurait exception si l'on avait $h = q$). De sorte qu'on aura la congruence
\[ \sum \beta B(q - 1, h', j', k') \equiv \sum \gamma B(q, h'', j'', k''). \tag{4} \]

D'ailleurs, comme $B(q, h'', j'', k'')$ est sur la frontière de $a^{h}_{j}$, il en sera de même de $a^{h''}_{j''}$; car si $B(q, h'', j'', k'')$ fait partie à la fois de $a^{h''}_{j''}$ et d'une variété $a^{h'}_{j'}$, faisant partie de la frontière de $a^{h}_{j}$; ou bien $a^{h''}_{j''}$ ne fait pas partie de $a^{h'}_{j'}$, et alors $B$ devrait faire partie de $a^{m}_{k}$, où $m < h''$, et nous avons vu que cela était impossible.

\origpage{306}

On a donc
\[ h'' < h. \]

Les congruences (2) et (4) donnent
\[ S \alpha B(q, h, j, k) \equiv \sum \gamma B(q, h'', j'', k''), \]
et, comme toutes les variétés qui figurent dans cette congruence, font partie de $a^{h}_{j}$, ou de sa frontière, comme, d'autre part, $a^{h}_{j}$ est simplement connexe, on aura l'homologie
\[ S \alpha B(q, h, j, k) \backsim \sum \gamma B(q, h'', j'', k''). \]

On peut, donc, remplacer, dans le premier membre de (1), l'ensemble des termes $S \alpha B(q, h, j, k)$ par l'ensemble des termes
\[ \sum \gamma B(q, h'', j'', k''). \]

Si on opère de même pour toutes les classes correspondantes à une même valeur de $h$, la plus grande de toutes, on aura remplacé le premier membre de (1) par
\[ \sum \alpha_{2} B(q, h_{2}, j_{2}, k_{2}), \]
où la plus grande valeur de $h_{2}$ sera plus petite que la plus grande valeur de $h$. On aura, d'ailleurs, l'homologie
\[ \sum \alpha B(q, h, j, k) \backsim \sum \alpha_{2} B(q, h_{2}, j_{2}, k_{2}). \]

En continuant de la sorte, on pourra diminuer encore la plus grande valeur de $h$. On ne sera arrêté que quand on aura partout $h = q$.

On peut donc, finalement, remplacer le premier membre de (1) par~:
\[ \sum \alpha_{0} B(q, q, j_{0}, k_{0}), \]
\origpage{307}
et on aura d'ailleurs~:
\[ \sum \alpha B(q, h, j, k) \backsim \sum \alpha_{0} B(q, q, j_{0}, k_{0}) \tag{5} \]
\[ \sum \alpha_{0} B(q, q, j_{0}, k_{0}) \equiv 0. \tag{6} \]

Cela posé, dans le premier membre de (6) prenons les congruences qui appartiennent à une classe déterminée $a^{h}_{j_{0}}$; soit~:
\[ S \alpha_{0} B(q, q, j_{0}, k_{0}). \]

Nous aurons~:
\[ S \alpha_{0} B(q, q, j_{0}, k_{0}) \equiv \sum \beta_{0} B(q - 1, h'_{0}, j'_{0}, k'_{0}). \tag{7} \]

Nous verrions, comme plus haut, que $a^{h'_{0}}_{j'_{0}}$ doit faire partie de la frontière de $a^{q}_{j_{0}}$, d'où $h'_{0} < q$ (et, comme $h'_{0} \geqq q - 1$, on aura $h'_{0} = q - 1$).

Soit alors~:
\[ a^{q}_{j_{0}} = \sum B(q, q, j_{0}, k_{0}) \tag{8} \]
l'équation (1, $q$, $j_{0}$), qui définit la subdivision de la variété $a^{q}_{j_{0}}$, et soient $B(q, q, j_{0}, 1)$ et $B(q, q, j_{0}, 2)$ deux variétés, figurant dans le second membre de (8); je dis qu'elles devront figurer dans le premier membre de (6) avec le même coefficient $\alpha_{0}$.

Supposons, d'abord, que ces deux variétés soient limitrophes; parmi les variétés à $q - 1$ dimensions qui leur serviront de frontière commune, il y en aura, au moins, une qui n'appartiendra pas à la frontière de $a^{q}_{j_{0}}$, qui fera, par conséquent, partie de la classe $a^{q}_{j_{0}}$.

Soit $B(q - 1, q, j_{0}, 1)$ cette variété~: elle n'appartiendra pas à aucune autre des variétés $B(q, q, j_{0}, k)$.

Soit alors
\[ \begin{aligned} B(q, q, j_{0}, 1) &\equiv \varepsilon b^{q-1}_{i} \\ B(q, q, j_{0}, 2) &\equiv \varepsilon b^{q-1}_{i} \\ B(q, q, j_{0}, k) &\equiv \varepsilon b^{q-1}_{i} \qquad (k > 2) \end{aligned} \tag{9} \]
\origpage{308}
les congruences (2, $q$, $q$, $j_{0}$, 1), (2, $q$, $q$, $j_{0}$, 2), (2, $q$, $q$, $j_{0}$, $k$), qui nous font connaitre les frontières des variétés $B(q, q, j_{0})$. Voyons avec quel coefficient $\varepsilon$ la variété $B(q - 1, q, j_{0}, 1)$ figurera dans ces congruences.

D'après ce que nous venons de voir, ce sera avec le coefficient $+1$ dans la première, avec le coefficient $-1$ dans la seconde, avec le coefficient 0 dans les autres.

Soient donc $\alpha_{1}$ et $\alpha_{2}$ les valeurs des coefficients $\alpha_{0}$, correspondantes aux deux variétés $B(q, q, j_{0}, 1)$ et $B(q, q, j_{0}, 2)$.

La combinaison des congruences (9) nous fournira une congruence
\[ S \alpha_{0} B(q, q, j_{0}, k_{0}) \equiv \varepsilon b^{q-1}_{i}, \tag{10} \]
qui devra être identique à (7), et le coefficient $\varepsilon$, avec lequel figurera $B(q - 1, q, j_{0}, 1)$ dans le second membre de (10), sera évidemment $\alpha_{1} - \alpha_{2}$. Mais $B(q - 1, q, j_{0}, 1)$ ne peut pas figurer dans le second membre de (7), puisque nous avons vu que dans ce second membre on doit avoir $h'_{0} = q - 1$. Donc on doit avoir $\alpha_{1} - \alpha_{2} = 0$.

Ainsi les deux variétés $B(q, q, j_{0}, 1)$ et $B(q, q, j_{0}, 2)$ devront avoir même coefficient $\alpha_{0}$, si elles sont limitrophes. Cela sera encore vrai, si elles ne le sont pas, parce que $a^{q}_{j_{0}}$ étant d'un seul tenant, on pourra passer d'une de ces variétés à l'autre par une suite d'autres variétés analogues, chacune d'elles étant limitrophe de celle qui la précède.

Donc le coefficient $\alpha_{0}$ est le même pour toutes nos variétés. D'où~:
\[ S \alpha_{0} B(q, q, j_{0}, k_{0}) = \alpha_{0} \sum B(q, q, j_{0}, k_{0}) = \alpha_{0} a^{q}_{j_{0}}. \]

La congruence (6) et l'homologie (5) peuvent donc s'écrire~:
\[ \sum \alpha B(q, h, j, k) \backsim \sum \alpha_{0} a^{q}_{j_{0}} \tag{5 \textit{bis}} \]
\[ \sum \alpha_{0} a^{q}_{j_{0}} \equiv 0. \tag{6 \textit{bis}} \]

Si un nombre quelconque de congruences de la forme (1) sont distinctes, c'est-à-dire si aucune combinaison linéaire de leurs premiers
\origpage{309}
membres n'est pas homologue à 0, je dis que les congruences (6~\emph{bis}) seront également distinctes, et réciproquement.

En effet la comparaison des relations (1), (5~\emph{bis}) et (6~\emph{bis}) montre que si l'on a
\[ \sum \alpha B(q, h, j, k) \backsim 0, \]
on aura également
\[ \sum \alpha_{0} a^{q}_{j_{0}} \backsim 0, \]
et réciproquement.

Il résulte de là que si les $a^{q}_{i}$ et les $b^{q}_{i}$ sont simplement connexes, les nombres de \emph{Betti} réduits sont les mêmes pour les deux polyèdres $V$ et $V'$.

Soit maintenant $W$ une variété quelconque, fermée, à $q$ dimensions, située sur $V$. On peut toujours construire un polyèdre $V'$, dérivé de $V$, au sens de la page 101 de l'\emph{Analysis situs}, et tel que $W$ soit une combinaison des $b^{q}_{i}$.

Nous devons donc conclure que~: si les $a^{q}_{i}$ sont simplement connexes, les nombres de \emph{Betti} réduits, relatifs au polyèdre $V$, sont identiques aux nombres de \emph{Betti} proprement dits, \emph{définis de la seconde manière}.

\subsection*{\S~VI. Retour sur les démonstrations du \S~III.}

Nous avons à revenir ici sur un point essentiel du raisonnement qui précède. J'ai dit plus haut qu'il n'y avait d'autre homologie que les homologies (9, $q$), obtenues au \S~3. Cela n'est pas évident, cela ne serait pas même toujours vrai, si nous ne supposions pas les $a^{q}_{i}$ simplement connexes.

Démontrons-le d'abord pour un polyèdre $P$, dans l'espace à 4 dimensions.

Considérons un certain nombre de variétés $v_{2}$ ou $a^{2}_{i}$, appartenant à ce polyèdre; je les appellerai ses faces, de même que les $a^{3}_{i}$, les $a^{1}_{i}$ et les $a^{0}_{i}$ de ce polyèdre $P$ pourront s'appeler ses cases, ses arêtes et ses sommets.

\origpage{310}

Supposons que l'on ait entre ces faces $a^{2}_{i}$ une homologie
\[ \sum a^{2}_{i} \backsim 0. \]

Cette homologie signifie qu'il existe une variété à 3 dimensions, $V$, faisant partie de $P$, et admettant $\displaystyle\sum a^{2}_{i}$ comme frontière complète.

Je dis que \emph{$V$ se compose d'un certain nombre de cases de $P$}.

Si, en effet, un point d'une case appartient à $V$, il en sera de même de tout autre point de cette case, car on peut aller du premier point au second, sans rencontrer aucune face et, par conséquent, sans rencontrer la frontière de $V$ et sans sortir de $V$.

Le théorème est donc évident en ce qui concerne les polyèdres de l'espace à 4 dimensions et les homologies entre les faces.

Soit maintenant une homologie entre les arêtes~:
\[ \sum b_{1} \backsim 0, \]
les $b_{1}$ étant un cetain\ednote{Printed ``cetain'' for ``certain''; an apparent misprint.} nombre d'arêtes $a^{1}_{i}$. Cela veut dire qu'il existe une variété à 2 dimensions, $V$, dont $\displaystyle\sum b_{1}$ est la frontière complète.

Je désignerai par $V(a^{k}_{i})$ l'ensemble des points communs à $V$ et à $a^{k}_{i}$.

Les $V(a^{3}_{i})$ seront des variétés à 2 dimensions, dont la frontière sera formée soit par quelques-unes des arêtes $b_{1}$, soit par les $V(a^{2}_{j})$, les $a^{2}_{j}$ étant les faces qui servent de frontière à la case $a^{3}_{i}$. On ne peut, en effet, sortir de $V(a^{3}_{i})$ qu'en sortant de $V$ par sa frontière, c'est-à-dire, en traversant une des $b_{1}$, ou qu'en sortant de $a^{3}_{i}$ par sa frontière, c'est-à-dire en traversant une face $a^{2}_{j}$, et, comme on reste sur $V$, en traversant une des lignes $V(a^{2}_{j})$.

La variété totale $V$ est formée de l'ensemble des $V(a^{3}_{i})$.

Considérons maintenant $V(a^{2}_{i})$; nous devons distinguer deux cas~:

1° ou bien aucune des arêtes $b_{1}$ n'appartient à $a^{2}_{i}$. Nous ne pourrons alors sortir de $V(a^{2}_{i})$ qu'en sortant de $a^{2}_{i}$, c'est-à-dire en traversant une des arêtes $a^{1}_{j}$; la frontière de $V(a^{2}_{i})$ est donc formée par les $V(a^{1}_{j})$.

\origpage{311}

2° ou bien une (ou plusieurs) arête $b_{1}$ fait partie de $a^{2}_{i}$; dans ce cas elle fera également partie de $V(a^{2}_{i})$; mais il pourra se faire que $V(a^{2}_{i})$ se compose, outre l'arête $b_{1}$, d'autres lignes; ces lignes auront pour frontières des points $V(a^{1}_{j})$, ou des points situés sur $b_{1}$. Ces points situés sur $b_{1}$, et où les autres lignes, dont se compose $V(a^{2}_{i})$, viennent se terminer sur l'arête $b_{1}$, seront ce que j'appellerai des \emph{points nodaux}.

Dans tous les cas $V(a^{2}_{i})$ sera une ligne ou un ensemble de lignes; si, en effet, $V(a^{2}_{i})$ était une surface, c'est que $a^{2}_{i}$, ou une portion de cette face, ferait partie de $V$. Mais j'ai le droit de déformer $V$, pourvu que je ne change pas sa frontière $\displaystyle\sum b_{1}$; je puis toujours, par une déformation infiniment petite, éviter qu'une région de $a^{2}_{i}$ fasse partie de $V$.

Pour la même raison je puis toujours supposer que $V(a^{1}_{i})$ se réduit à un ou plusieurs points, sauf si $a^{1}_{i}$ est l'une des arêtes $b_{1}$, auquel cas $V(a^{1}_{i})$ sera cette arête elle-même.

Cela posé, je puis déformer $V$~:

1° de manière que tous les $V(a^{1}_{i})$ [autres que $V(b_{1})$] soient des sommets. Soit $a^{0}_{j}$ un sommet de $a^{1}_{i}$. Soit $M$ l'un des points dont se compose $V(a^{1}_{i})$; autour du point $M$ et sur $V$ décrivons une petite courbe fermée $C$. Soit $K$ l'aire infiniment petite découpée sur $V$ par cette courbe $C$. Construisons une sorte de manchon, infiniment délié, entourant l'arête $a^{1}_{i}$ et passant par $C$. Par le sommet $a^{0}_{j}$ je mène une surface quelconque $S$; elle viendra découper sur le manchon une courbe fermée très petite $C'$. Soit $K'$ la portion de la surface $S$ limitée par $C$\ednote{Printed $C$; since $K'$ is the portion of $S$ cut out by the small curve $C'$, $C'$ appears to be intended.}; soit $H$ la surface du manchon comprise entre $C$ et $C'$. On figurera ainsi une sorte de tambour, dont $H$ sera la surface latérale, $K$ et $K'$ les deux bases.

Considérons alors la variété
\[ V' = V - K + H + K'. \]

Cette variété aura même frontière que $V$; mais elle ne coupera plus $a^{1}_{i}$ en $M$, puisqu'on a supprimé la portion $K$ de $V$, où se trouvait ce point $M$. En revanche, $H$ ne coupera pas l'arête $a^{1}_{i}$, et $K'$ coupera cette arête en $a^{0}_{j}$.

\origpage{312}

En opérant de même sur tous les points d'intersection de $V$ et de $a^{1}_{i}$, on amènerait tous ces points à coincider avec $a^{0}_{j}$.

2° de manière que tous les points nodaux soient des sommets.

Soit, en effet, $a^{2}_{i}$ une face passant par l'arête $b_{1}$; l'intersection de $V$ et $a^{2}_{i}$ comprendra, outre $b_{1}$, d'autres lignes; soit $c$ l'une de ces lignes, venant se terminer sur $b_{1}$ en un point nodal $D$. Soient $a^{0}_{j}$ et $a^{0}_{k}$ les deux sommets de $b_{1}$. Par $b_{1}$ je fais passer une surface $S$, faisant partie de $P$ et ne coupant pas $a^{2}_{i}$. Comme $a^{0}_{j}$ et $a^{0}_{k}$ sont sur la frontière de $V$, je joins ces deux points par une ligne $L$, située sur $V$ et s'écartant peu de $b_{1}$. Cette ligne s'écartant peu de $b_{1}$, je puis mener par $L$ une autre surface $S'$, qui ne passera pas par $b_{1}$, mais qui coupera $S$ suivant une ligne $L'$, très peu différente de $b_{1}$. Ces trois lignes $L$, $b_{1}$ et $L'$ auront mêmes extrémités $a^{0}_{j}$ et $a^{0}_{k}$.

Soit $V_{1}$ la portion de $V$, comprise entre $L$ et $b_{1}$; soit $S_{1}$ la portion de $S$, comprise entre $L'$ et $b_{1}$, et $S_{1}$\ednote{Printed $S_{1}$ here and in the display below, the same letter as the portion of $S$; a distinct name such as $S'_{1}$ for the portion of $S'$ appears to be intended.} la portion de $S'$, comprise entre $L$ et $L'$.

Je remplace $V$ par
\[ V' = V - V_{1} + S_{1} + S_{1}. \]
$V'$ a mêmes frontières que $V$, mais $V'(a^{2}_{i})$ ne présente plus de points nodaux, en dehors de $a^{0}_{j}$ et $a^{0}_{k}$; car si une ligne analogue à $c$ venait aboutir à un point nodal, situé entre $a^{0}_{j}$ et $a^{0}_{k}$, la portion de cette ligne $c$, voisine de ce point nodal, devrait se trouver sur $S_{1}$, ce qui est impossible, puisque $S$ ne coupe pas $a^{2}_{i}$.

En résumé~: nous pouvons toujours supposer que les $V(a^{2}_{i})$ sont des lignes, dont les extrémités sont des sommets de $a^{2}_{i}$.

Soit alors $L$ une des lignes, dont se compose $V(a^{2}_{i})$, ayant pour extrémités deux sommets $a^{0}_{j}$ et $a^{0}_{k}$ de $a^{2}_{i}$. On peut aller de $a^{0}_{j}$ à $a^{0}_{k}$, en suivant le périmètre de $a^{2}_{i}$; soit $\displaystyle\sum a^{1}_{m}$ l'ensemble des arêtes de $a^{2}_{i}$, comprises entre $a^{0}_{j}$ et $a^{0}_{k}$. Comme la face $a^{2}_{i}$ est supposée simplement connexe, la ligne $L$ la divisera en deux parties. Soit $Q$ l'une de ces parties, comprise entre $L$ et $\displaystyle\sum a^{1}_{m}$.

Soient $a^{3}_{p}$ et $a^{3}_{q}$ les deux cases séparées par $a^{2}_{i}$. Par les arêtes $\displaystyle\sum a^{1}_{m}$ je fais passer une surface $S$, peu différente de la face $a^{2}_{i}$ et située toute entière dans la case $a^{3}_{p}$; par les mêmes arêtes je fais passer une seconde
\origpage{313}
surface $S'$, peu différente de $a^{2}_{i}$ et située dans la case $a^{3}_{q}$; ces deux surfaces $S$ et $S'$ couperont $V$, suivant deux lignes $L_{1}$ et $L'_{1}$, peu différentes de $L$, et ayant pour extrémités $a^{0}_{j}$ et $a^{0}_{k}$. Soit $S_{1}$ la portion de $S$ comprise entre $\displaystyle\sum a^{1}_{m}$ et $L_{1}$; soit $S_{1}$\ednote{Printed $S_{1}$ here and in the display below, the same letter as the portion of $S$; $S'_{1}$, the portion of $S'$, appears to be intended.} la portion de $S'$ comprise entre $\displaystyle\sum a^{1}_{m}$ et $L'_{1}$; soit $V_{1}$ la portion de $V$ comprise entre $L_{1}$ et $L'_{1}$; c'est sur $V_{1}$ que se trouvera $L$.

Soit maintenant
\[ V' = V - V_{1} + S_{1} + S_{1}. \]
$V'$ a mêmes frontières que $V$; $V'$ ne passe plus par $L$, mais en revanche passe par les arêtes $\displaystyle\sum a^{1}_{m}$.

En opérant de la même manière pour toutes les lignes telles que $L$, on voit qu'on peut toujours supposer que tous les $V(a^{2}_{i})$ se réduisent à des combinaisons d'arêtes.

Comme les frontières de $V(a^{3}_{i})$ sont ou des $b_{1}$ ou des $V(a^{2}_{j})$, on voit que les frontières des $V(a^{3}_{i})$ sont des combinaisons d'arêtes de $P$, et, bien entendu, toutes ces arêtes doivent appartenir à $a^{3}_{i}$. Ainsi donc $V(a^{3}_{i})$ est une surface simplement ou multiplement connexe, limitée par une ou plusieurs lignes fermées, qui, elles-mêmes, sont des combinaisons d'arêtes de $a^{3}_{i}$.

Comme la case $a^{3}_{i}$ est simplement connexe, ces lignes fermées subdiviseront la surface de cette case en un certain nombre de régions, et comme ces lignes fermées sont des combinaisons d'arêtes de $a^{3}_{i}$, ces régions seront des combinaisons des faces de $a^{3}_{i}$.

On pourra toujours trouver une combinaison de ces régions, qui aura mêmes frontières que $V(a^{3}_{i})$. Supposons, par exemple, que la frontière de $V(a^{3}_{i})$ se compose de trois lignes fermées $L, L_{1}, L_{2}$; la ligne $L$ divisera la surface de $a^{3}_{i}$ en deux régions $R$ et $R'$; les lignes $L_{1}$ et $L_{2}$ diviseront de même cette surface en deux régions $R_{1}$ et $R'_{1}$, ou $R_{2}$ et $R'_{2}$. Je suppose qu'en parcourant $L$ dans un certain sens, on ait à sa gauche $V(a^{3}_{i})$, et $R$ et $R'$ à sa droite; je suppose de même qu'en parcourant $L_{1}$ et $L_{2}$ dans un sens convenable, on ait à sa gauche $V(a^{3}_{i})$ et $R_{1}$, ou $V(a^{3}_{i})$ et $R_{2}$.

Alors la variété $R + R_{1} + R_{2}$ aura même frontière que $V(a^{3}_{i})$; on pourra donc remplacer $V(a^{3}_{i})$ par $R + R_{1} + R_{2}$.

\origpage{314}

En opérant de la même manière sur tous les $V(a^{3}_{i})$, on aura remplacé $V$ par une autre variété, qui aura même frontière $\displaystyle\sum b_{1}$, et qui sera une combinaison de faces de $P$.

Le théorème est donc démontré en ce qui concerne les polyèdres de l'espace à 4 dimensions et les homologies entre les arêtes.

On le démontrerait de même pour un polyèdre quelconque.

\subsection*{\S~VII. Polyèdre réciproque.}

Soit $P$ un polyèdre dans l'espace à 4 dimensions; ce polyèdre sera subdivisé en un certain nombre de variétés $v_{3}$, que j'appellerai ses \emph{cases}, et que je désignerai par $a^{3}_{i}$. Ces cases seront separées les unes des autres par des variétés $v_{2}$ ou $a^{2}_{i}$, que j'appellerai les \emph{faces}; ces faces auront pour frontières des variétés $v_{1}$, ou $a^{1}_{i}$, que j'appellerai les \emph{arêtes}, et les extrémités des arêtes seront des points $v_{0}$ ou $a^{0}_{i}$, que j'appellerai les \emph{sommets}.

Je supposerai, bien entendu, que les cases et les faces sont simplement connexes.

Marquons, à l'intérieur de chaque case $a^{3}_{i}$, un point $P(a^{3}_{i})$; à l'intérieur de chaque face $a^{2}_{i}$, un point $P(a^{2}_{i})$; sur chaque arête $a^{1}_{i}$, un point $P(a^{1}_{i})$; chaque arête se trouvera ainsi partagée en deux parties par le point $P(a^{1}_{i})$.

Joignons par des lignes le point $P(a^{2}_{i})$ à chacun des sommets de la face $a^{2}_{i}$ et à chacun des points $P(a^{1}_{j})$, correspondant aux diverses arêtes $a^{1}_{j}$ de la face $a^{2}_{i}$. Toutes ces lignes devront être tracées sur la face $a^{2}_{i}$. Cette face sera ainsi partagée en triangles, et le nombre de ces triangles sera double du nombre des arêtes de $a^{2}_{i}$. Nous ferons de même pour toutes les autres faces.

Considérons maintenant une case $a^{3}_{i}$; décomposons en triangles $T$ toutes les faces $a^{2}_{j}$ de cette case, ainsi que nous venons de le dire. Construisons des triangles curvilignes, ayant pour sommet commun le point $P(a^{3}_{i})$ et pour bases les différents côtés des différents triangles $T$. La case $a^{3}_{i}$ sera ainsi décomposée en tétraèdres, ayant $P(a^{3}_{i})$ pour sommet commun et pour bases les différents triangles $T$.

\origpage{315}

Nous distinguerons six sortes de lignes (qui seront les arêtes de nos tétraèdres)~:

celles de la 1\textsuperscript{ère} sorte joindront un sommet $a^{0}_{i}$ à un point $P(a^{1}_{j})$; chaque arête sera ainsi formée de deux lignes de la 1\textsuperscript{ère} sorte;

celles de la 2\textsuperscript{de} sorte joindront un point $P(a^{3}_{i})$ à un point $P(a^{2}_{j})$;

» » 3\textsuperscript{e} » » » $P(a^{2}_{i})$ » sommet $a^{0}_{j}$;

» » 4\textsuperscript{e} » » » $P(a^{1}_{i})$ » point $P(a^{2}_{j})$;

» » 5\textsuperscript{e} » » » $P(a^{3}_{i})$ » » $P(a^{1}_{j})$;

» » 6\textsuperscript{e} » » » » » sommet $a^{0}_{j}$.

Les lignes de la 2\textsuperscript{de} sorte peuvent s'accoupler deux à deux de deux manières~:

1° Ce que j'appellerai la ligne $b^{1}_{i}$ sera formée de deux lignes de la 2\textsuperscript{de} sorte, joignant un même point $P(a^{2}_{i})$ à deux points $P(a^{3}_{j})$ et $P(a^{3}_{k})$, correspondant aux deux cases $a^{3}_{j}$ et $a^{3}_{k}$ séparées par la face $a^{2}_{i}$. Il y aura donc autant de lignes $b^{1}_{i}$ que de faces $a^{2}_{i}$.

2° Ce que j'appellerai une ligne $c$ sera formée de deux lignes de la 2\textsuperscript{de} sorte, joignant un même point $P(a^{3}_{i})$ à deux points $P(a^{2}_{j})$ et $P(a^{2}_{k})$, correspondant à deux faces $a^{2}_{j}$ et $a^{2}_{k}$ de la case $a^{3}_{i}$.

Il nous faut définir des surfaces que j'appellerai les surfaces $b^{2}_{i}$.

Soit une arête quelconque $a^{1}_{i}$ et le point $P(a^{1}_{i})$. Supposons que les faces qui passent par $a^{1}_{i}$, soient successivement~:
\[ a^{2}_{1}, \qquad a^{2}_{2}, \qquad \ldots, \qquad a^{2}_{q}, \]
et que les cases, auxquelles appartient $a^{1}_{i}$, soient successivement~:
\[ a^{3}_{1}, \qquad a^{3}_{2}, \qquad \ldots, \qquad a^{3}_{q}, \]
de telle façon que $a^{2}_{1}$ sépare $a^{3}_{1}$ de $a^{3}_{2}$, $a^{2}_{2}$ sépare $a^{3}_{2}$ de $a^{3}_{3}$, \ldots, et, qu'enfin, $a^{2}_{q}$ sépare $a^{3}_{q}$ de $a^{3}_{1}$. Convenons, pour plus de symétrie, de désigner indifféremment la case $a^{3}_{1}$ par $a^{3}_{1}$ ou $a^{3}_{q+1}$.

Décomposons chaque case en tétraèdres, et envisageons en particulier les tétraèdres qui admettent pour sommet le point $P(a^{1}_{i})$. Considérons les $2q$ triangles curvilignes~:
\[ P(a^{1}_{i}) P(a^{3}_{k}) P(a^{2}_{k}), \qquad P(a^{1}_{i}) P(a^{3}_{k+1}) P(a^{2}_{k}). \qquad (k = 1, 2, \ldots, q) \]

\origpage{316}

L'ensemble de ces $2q$ triangles formera un certain polygone que j'appellerai $b^{2}_{i}$, et qui aura pour frontière l'ensemble des lignes
\[ b^{1}_{1}, \qquad b^{1}_{2}, \qquad \ldots, \qquad b^{1}_{q}. \]

Définissons maintenant les volumes $b^{3}_{i}$; le volume $b^{3}_{i}$ sera l'ensemble des tétraèdres qui admettent pour sommet le point $a^{0}_{i}$; ce volume sera un polyèdre à 3 dimensions, simplement connexe, qui aura pour frontière l'ensemble des surfaces $b^{2}_{k}$, correspondant aux arêtes $a^{1}_{k}$, qui aboutissent au point $a^{0}_{i}$.

La juxtaposition des volumes $b^{3}_{i}$ constituera un nouveau polyèdre $P'$, que j'appellerai le \emph{polyèdre réciproque de $P$}, et qui aura pour cases le\ednote{Printed ``le'' for ``les''; an apparent misprint.} $b^{3}_{i}$, pour faces les $b^{2}_{i}$, pour arêtes les $b^{1}_{i}$, pour sommets les points $b^{0}_{i} = P(a^{3}_{i})$.

A chaque case $b^{3}_{i}$ de $P'$ correspondra un sommet $a^{0}_{i}$ de $P$,

» » face $b^{2}_{i}$ » » » une arête $a^{1}_{i}$ » »

» » arête $b^{1}_{i}$ » » » » face $a^{2}_{i}$ » »

» » sommet $b^{0}_{i}$ » » » » case $a^{3}_{i}$ » ».

De plus, au sens du \S~II, il y aura la même relation, par exemple, entre l'arête $b^{1}_{i}$ et la face $b^{2}_{j}$, qu'entre la face $a^{2}_{i}$ et l'arête $a^{1}_{j}$.

Si donc les congruences caractéristiques du polyèdre $P$ s'écrivent~:
\[ a^{3}_{i} \equiv \sum_{j} \varepsilon^{3}_{i,j} a^{2}_{i}, \qquad a^{2}_{i} \equiv \sum_{j} \varepsilon^{2}_{i,j} a^{1}_{i}, \qquad a^{1}_{i} \equiv \sum_{j} \varepsilon^{1}_{i,j} a^{0}_{j}, \]
\ednote{The first two sums are printed with $a^{2}_{i}$ and $a^{1}_{i}$ under the sum over $j$; $a^{2}_{j}$ and $a^{1}_{j}$, as in the third, appear to be intended.}
celles du polyèdre $P'$ s'écriront~:
\[ b^{3}_{i} \equiv \sum_{j} \varepsilon^{1}_{i,j} b^{2}_{j}, \qquad b^{2}_{i} \equiv \sum_{j} \varepsilon^{2}_{i,j} b^{1}_{j}, \qquad b^{1}_{i} \equiv \sum_{j} \varepsilon^{3}_{i,j} b^{0}_{j}. \]

Considérons maintenant une ligne $c$, formée de deux lignes de la seconde sorte, joignant un même point $P(a^{3}_{i})$ à deux points $P(a^{2}_{j})$ et $P(a^{2}_{k})$.

Soient $a^{0}_{m}$ et $a^{0}_{p}$ deux sommets, appartenant respectivement tous deux à la case $a^{3}_{i}$. Soient $d$ et $d'$ les deux lignes de la 3\textsuperscript{ème} sorte qui joignent respectivement $P(a^{2}_{j})$ à $a^{0}_{m}$, et $P(a^{2}_{k})$ à $a^{0}_{p}$.

Comme $a^{0}_{m}$ et $a^{0}_{p}$ appartiennent à une même case $a^{3}_{i}$, on pourra aller de l'un de ces sommets à l'autre, en suivant une ligne brisée formée d'arêtes $a^{1}_{q}$ appartenant à $a^{3}_{i}$.

\origpage{317}

Soit $\displaystyle\sum a^{1}_{q}$ cette ligne brisée, dont les extrémités sont $a^{0}_{m}$ et $a^{0}_{p}$; l'ensemble des lignes $c - d - \displaystyle\sum a^{1}_{q} + d'$ sera une ligne fermée, ce que j'exprimerai par la congruence~:
\[ c \equiv d + \sum a^{1}_{q} - d'. \]

Comme $a^{3}_{i}$ est simplement connexe, cette ligne fermée sera la frontière d'une variété à 2 dimensions, intérieure à $a^{3}_{i}$, ce que j'exprimerai par l'homologie~:
\[ c \backsim d + \sum a^{1}_{q} - d'. \]

Réciproquement~: soit $\displaystyle\sum a^{1}_{q}$ une ligne brisée, formée d'arêtes appartenant toutes à $a^{3}_{i}$, et dont les extrémités sont les sommets $a^{0}_{m}$ et $a^{0}_{p}$; ces deux sommets appartiendront respectivement à deux faces $a^{2}_{j}$ et $a^{2}_{k}$, faisant toutes deux partie de $a^{3}_{i}$. Soient les trois lignes~:
\[ c = P(a^{2}_{j}) P(a^{3}_{i}) + P(a^{3}_{i}) P(a^{2}_{k}), \qquad d = P(a^{2}_{j}) a^{0}_{m}, \qquad d' = P(a^{2}_{k}) a^{0}_{p}. \]

On aura encore
\[ c \backsim d + \sum a^{1}_{q} - d'. \]

Soit maintenant $a^{0}_{i}$ un sommet appartenant à deux faces $a^{2}_{j}$ et $a^{2}_{k}$. Soient les deux lignes de la 3\textsuperscript{ème} sorte~:
\[ d_{j} = P(a^{2}_{j}) a^{0}_{i}, \qquad d_{k} = P(a^{2}_{k}) a^{0}_{i}. \]

Nous pouvons tracer une ligne $L$, s'écartant infiniment peu du sommet $a^{0}_{i}$, et allant d'un point de $a^{2}_{j}$ à un point de $a^{2}_{k}$.

Supposons, pour fixer les idées, que cette ligne traverse trois cases et qu'elle rencontre successivement la face $a^{2}_{j}$, la case $a^{3}_{j}$, la face $a^{2}_{m}$, la case $a^{3}_{m}$, la face $a^{2}_{p}$, la case $a^{3}_{p}$, et enfin la face $a^{2}_{k}$.

Construisons les trois lignes $c$~:
\[ \begin{aligned} c_{j} &= P(a^{2}_{j}) P(a^{3}_{j}) + P(a^{3}_{j}) P(a^{2}_{m}), \\ c_{m} &= P(a^{2}_{m}) P(a^{3}_{m}) + P(a^{3}_{m}) P(a^{2}_{p}), \\ c_{p} &= P(a^{2}_{p}) P(a^{3}_{p}) + P(a^{3}_{p}) P(a^{2}_{k}), \end{aligned} \]
\origpage{318}
et les deux lignes de la 3\textsuperscript{ème} sorte~:
\[ d_{m} = P(a^{2}_{m}) a^{0}_{i}, \qquad d_{p} = P(a^{2}_{p}) a^{0}_{i}. \]

On aura~:
\[ c_{j} \equiv d_{j} - d_{m}, \qquad c_{m} \equiv d_{m} - d_{p}, \qquad c_{p} \equiv d_{p} - d_{k}; \]
et comme les trois cases $a^{3}_{j}, a^{3}_{m}, a^{3}_{p}$ sont simplement connexes~:
\[ c_{j} \backsim d_{j} - d_{m}, \qquad c_{m} \backsim d_{m} - d_{p}, \qquad c_{p} \backsim d_{p} - d_{k}, \]
et enfin~:
\[ c_{j} + c_{m} + c_{p} \backsim d_{j} - d_{k}. \]

On peut donc toujours trouver une ligne brisée, formée de lignes $c$ et homologue à $d_{j} - d_{k}$, $d_{j}$ et $d_{k}$ étant des lignes de la 3\textsuperscript{ème} sorte, aboutissant à un même sommet.

Cela posé, soit
\[ \sum b^{1}_{i} \equiv 0 \tag{1} \]
une congruence entre les arêtes $b^{1}_{i}$ du polyèdre $P'$.

La ligne brisée $\displaystyle\sum b^{1}_{i}$ est évidemment formée d'un nombre pair de lignes de la 2\textsuperscript{ème} sorte, et en parcourant cette ligne brisée, on rencontrera successivement $q$ faces
\[ a^{2}_{1}, \qquad a^{2}_{2}, \qquad \ldots, \qquad a^{2}_{q}, \]
pour revenir à la face $a^{2}_{1}$, que je désignerai également par $a^{2}_{q+1}$; et on rencontrera $q$ cases
\[ a^{3}_{1}, \qquad a^{3}_{2}, \qquad \ldots, \qquad a^{3}_{p}, \] \ednote{Printed $a^{3}_{p}$ as the last term; $a^{3}_{q}$ is intended, as the text speaks of $q$ cases and continues with $a^{3}_{q+1}$.}
pour revenir à la case $a^{3}_{1}$, que je désignerai également par $a^{3}_{q+1}$, de sorte que la face $a^{2}_{k}$ séparera la case $a^{3}_{k}$ de la case $a^{3}_{k+1}$.

Notre congruence s'écrit alors~:
\[ \sum [P(a^{3}_{k}) P(a^{2}_{k}) + P(a^{2}_{k}) P(a^{3}_{k+1})] \equiv 0, \]
\origpage{319}
ou, ce qui revient au même,
\[ \sum [P(a^{2}_{k-1}) P(a^{3}_{k}) + P(a^{3}_{k}) P(a^{2}_{k})] = 0. \]

Soit alors $a^{0}_{k}$ un sommet de la face $a^{2}_{k}$, appartenant, par conséquent, à la fois aux cases $a^{3}_{k}$ et $a^{3}_{k+1}$.

Soit $d_{k}$ la ligne de la 3\textsuperscript{ème} sorte $P(a^{2}_{k}) a^{0}_{k}$; nous venons de voir qu'il existe une ligne brisée $A_{k}$, formée d'arêtes appartenant à la case $a^{3}_{k}$, et telle que l'on ait l'homologie~:
\[ P(a^{2}_{k-1}) P(a^{3}_{k}) + P(a^{3}_{k}) P(a^{2}_{k}) \backsim d_{k-1} + A_{k} - d_{k}. \]

En additionnant toutes ces homologies, le premier membre se réduit à~:
\[ \sum [P(a^{2}_{k-1}) P(a^{3}_{k}) + P(a^{3}_{k}) P(a^{2}_{k})] = \sum b^{1}_{i}; \]
les lignes de la 3\textsuperscript{ème} sorte $d_{k}$ disparaissent, et il reste~:
\[ \sum b^{1}_{i} \backsim \sum A_{k}, \]
et par conséquent
\[ \sum b^{1}_{i} \equiv \sum A_{k} \equiv 0. \]

Donc, à toute congruence $\displaystyle\sum b^{1}_{i} \equiv 0$ entre les arêtes de $P'$, correspond une congruence $\displaystyle\sum A_{k} \equiv 0$ entre les arêtes de $P$, et telle que l'on ait~:
\[ \sum b^{1}_{i} \backsim \sum A_{k}. \]
Si donc on a $\displaystyle\sum b'_{i} \backsim 0$\ednote{Printed with a prime-like mark in place of the superscript 1, as $b'_{i}$; the context requires $b^{1}_{i}$.}, on aura $\displaystyle\sum A_{k} \backsim 0$, et réciproquement.

Soit maintenant
\[ \sum A_{k} \equiv 0 \tag{2} \]
une congruence entre les arêtes de $P$; supposons que $A_{k}$ soit une ligne brisée formée d'arêtes appartenant à la case $a^{3}_{k}$.

Le premier membre de la congruence (2) se composera de $q$
\origpage{320}
semblables lignes brisées~:
\[ A_{1}, \qquad A_{2}, \qquad \ldots, \qquad A_{q}, \]
et je désignerai indifféremment $A_{1}$ par $A_{1}$ ou $A_{q+1}$, et $A_{q}$ par $A_{0}$ ou $A_{q}$.

Soient $a^{0}_{k-1}$ et $a^{0}_{k}$ les deux extrémités de la ligne $A_{k}$; le sommet $a^{0}_{k}$ appartiendra à la fois aux cases $a^{3}_{k}$ et $a^{3}_{k+1}$; soit $a'^{2}_{k}$ la face de $a^{3}_{k}$, et $a^{2}_{k+1}$ la face de $a^{3}_{k+1}$, auxquelles appartient $a^{0}_{k}$.

Soient les lignes de la 3\textsuperscript{ème} sorte
\[ d'_{k} = P(a'^{2}_{k}) a^{0}_{k}, \qquad d_{k+1} = P(a^{2}_{k+1}) a^{0}_{k}, \]
et d'autre parte\ednote{Printed ``parte'' for ``part''; an apparent misprint.}~:
\[ c_{k} = P(a^{2}_{k}) P(a^{3}_{k}) + P(a^{3}_{k}) P(a'^{2}_{k}). \]

Nous avons vu que
\[ A_{k} \backsim - d_{k} + c_{k} + d'_{k}. \]

D'autre part, les lignes $d_{k+1}$ et $d'_{k}$ aboutissent à un même sommet $a^{0}_{k}$; nous avons vu également que l'on peut trouver une combinaison $C_{k}$ de lignes $c$, telle que l'on ait
\[ C_{k} \backsim d'_{k} - d_{k-1}. \] \ednote{Printed $d_{k-1}$; the preceding sentence names $d_{k+1}$ and $d'_{k}$ as the lines ending at $a^{0}_{k}$, so $d_{k+1}$ may be intended.}

En additionnant toutes ces homologies, je trouve~:
\[ \sum A_{k} \backsim \sum c_{k} + \sum C_{k}, \]
et par conséquent~:
\[ \sum c_{k} + \sum C_{k} \equiv 0. \]

Le premier membre de cette dernière congruence est une combinaison de lignes $c$, ou, ce qui revient au même, une combinaison d'arêtes $b^{1}_{i}$ du polyèdre $P'$, de sorte que je puis poser~:
\[ \sum c_{k} + \sum C_{k} = \sum b^{1}_{i}, \]
\origpage{321}
d'où
\[ \sum A_{k} \backsim \sum b^{1}_{i}. \]

En résumé~: à toute congruence entre les arêtes de $P$, correspond une congruence entre celles de $P'$, et réciproquement, et la condition nécessaire et suffisante pour que le premier membre de l'une des congruences soit homologue à zéro, c'est que l'autre le soit.

En d'autres termes, le nombre des congruences distinctes entre les arêtes est le même pour $P$ et $P'$, en ne considérant pas des congruences comme distinctes, quand une combinaison linéaire des premiers membres de ces congruences est homologue à zéro.

En d'autres termes encore, le nombre de \emph{Betti} réduit, relatif aux arêtes de $P$, est égal au nombre de \emph{Betti} réduit, relatif aux arêtes de $P'$.

On pourrait arriver au même résultat, en remarquant que l'on peut construire un polyèdre qui serait, à la fois, dérivé du polyèdre $P$ et dérivé du polyèdre réciproque $P'$, et en appliquant le théorème du \S~5.

Nous verrons plus loin, au \S~10, que cette proposition peut être présentée sous une autre forme.

D'autre part, cela peut permettre, plus simplement qu'au \S~5, de démontrer que les nombres de \emph{Betti} réduits sont égaux aux nombres de \emph{Betti} proprement dits.

En effet la définition du polyèdre $P'$ comporte un certain arbitraire~: ses sommets $b^{0}_{i}$ ne sont assujettis qu'à être intérieurs aux cases $a^{3}_{i}$ de $P$. Dans ces conditions, on peut évidemment choisir toujours le polyèdre $P'$ de façon qu'une ligne fermée \emph{quelconque} soit une combinaison des $b^{1}_{i}$.

\subsection*{\S~VIII. Démonstration du théorème fondamental.}

Soit $N_{1}$ le nombre des arêtes de notre polyèdre $P$, $N_{2}$ le nombre des faces, $N_{3}$ celui des cases. Formons un tableau d'après les règles suivantes.

\origpage{322}

Le tableau aura $N_{2} + N_{3}$ colonnes, $N_{2}$ dites de la 1\textsuperscript{ère} sorte et $N_{3}$ dites de la 2\textsuperscript{de}; il aura $N_{2} + N_{1}$ lignes, $N_{2}$ de la 1\textsuperscript{ère} et $N_{1}$ de la 2\textsuperscript{de} sorte. Voici quels seront les éléments du tableau~:

1° Pour l'élément de la $i^{\text{ème}}$ ligne de la 1\textsuperscript{ère} sorte et la $j^{\text{ème}}$ colonne de la 1\textsuperscript{ère} sorte, j'écrirai 1, si $i = j$, et 0, si $i \neq j$.

2° Les éléments appartenant à une ligne de la 2\textsuperscript{de} sorte et à une colonne de la 2\textsuperscript{de} sorte, seront tous nuls.

3° L'élément de la $i^{\text{ème}}$ colonne de la 1\textsuperscript{ère} sorte et de la $j^{\text{ème}}$ ligne de la 2\textsuperscript{de} sorte, sera $\varepsilon^{2}_{i,j}$, $\varepsilon^{2}_{i,j}$ étant le nombre qui nous fait connaître la relation entre la face $a^{2}_{i}$ et l'arête $a^{1}_{j}$.

4° L'élément de la $i^{\text{ème}}$ ligne de la 1\textsuperscript{ère} sorte et de la $j^{\text{ème}}$ colonne de la 2\textsuperscript{de} sorte, sera $\varepsilon^{3}_{i,j}$, c'est-à-dire le nombre qui fait connaître la relation entre la case $a^{3}_{j}$ et la face $a^{2}_{i}$.

Notre tableau, s'il y a par exemple deux cases, quatre faces et trois arêtes, présentera un aspect tel que celui-ci~:
\[ \begin{matrix} 1 & 0 & 0 & 0 & \varepsilon & \varepsilon \\[1ex] 0 & 1 & 0 & 0 & \varepsilon & \varepsilon \\[1ex] 0 & 0 & 1 & 0 & \varepsilon & \varepsilon \\[1ex] 0 & 0 & 0 & 1 & \varepsilon & \varepsilon \\[1ex] \varepsilon & \varepsilon & \varepsilon & \varepsilon & 0 & 0 \\[1ex] \varepsilon & \varepsilon & \varepsilon & \varepsilon & 0 & 0 \\[1ex] \varepsilon & \varepsilon & \varepsilon & \varepsilon & 0 & 0 \end{matrix} \tag{1} \]

Je n'ai pas écrit les indices des nombres $\varepsilon$ pour simplifier.

Voici maintenant les opérations que je regarde comme permises sur ce tableau.

1° Ajouter une colonne à une autre \emph{de même sorte}, ou l'en retrancher.

2° Ajouter une ligne à une autre \emph{de même sorte}, ou l'en retrancher.

3° Permuter deux colonnes de même sorte, en changeant tous les signes de l'une d'elles.

\origpage{323}

4° Permuter deux lignes de même sorte, en changeant tous les signes de l'une d'elles.

Toutes ces transformations, pour lesquelles les éléments du tableau restent entiers, s'appelleront les \emph{transformations arithmétiques} du tableau.

On peut s'en servir pour simplifier la partie du tableau qui appartient aux colonnes de la 1\textsuperscript{ère} sorte et aux lignes de la 2\textsuperscript{de} sorte, et celle qui appartient aux colonnes de la 2\textsuperscript{ème} sorte et aux lignes de la 1\textsuperscript{ère} sorte.

Voici jusqu'où l'on peut pousser la simplification, d'après des théorèmes bien connus d'arithmétique; quand la réduction sera terminée~:

L'élément de la $i^{\text{ème}}$ colonne de la 1\textsuperscript{ère} sorte et de la $j^{\text{ème}}$ ligne de la 2\textsuperscript{de} sorte~:

1° Sera nul, si $i > j$.

2° Sera égal à un entier $H_{i}$, qui pourra être nul, si $i = j$.

3° Sera encore nul, si $i < j$ et si $H_{i}$ est premier avec $H_{j}$.

4° Enfin sera nul, si $j > N_{2}$.

Il en sera de même de l'élément de la $i^{\text{ème}}$ ligne de la 1\textsuperscript{ère} sorte et de la $j^{\text{ème}}$ colonne de la 2\textsuperscript{de} sorte.

La réduction peut être poussée encore plus loin, si l'on autorise une cinquième opération~: multiplier tous les éléments d'une ligne ou d'une colonne par un même nombre entier ou non, différent de zéro.

Les transformations correspondantes s'appelleront les \emph{transformations algébriques} du tableau.

On peut alors supposer que l'élément de la $i^{\text{ème}}$ colonne de la 1\textsuperscript{ère} sorte et de la $j^{\text{ème}}$ ligne de la 2\textsuperscript{de} sorte (de même que l'élément de la $i^{\text{ème}}$ ligne de la 1\textsuperscript{ère} sorte et de la $j^{\text{ème}}$ colonne de la 2\textsuperscript{de} sorte) est nul, si $i \neq j$. Si $i = j$, il peut être égal à 0 ou à 1. S'il en est ainsi, je dirai que le tableau est \emph{réduit}.

Après la cinquième opération, les éléments qui appartiennent aux lignes et aux colonnes de la 1\textsuperscript{ère} sorte pourront ne pas rester entiers; de plus, le déterminant formé par ces lignes et ces colonnes pourra ne pas rester égal à 1, mais il restera différent de zéro.

Le tableau (1) est relatif aux faces du polyèdre $P$ et à leurs relations avec les cases et les arêtes. Nous pourrions en dresser un, tout pareil, relatif aux arêtes du polyèdre $P$ et à leurs relations avec les faces et les sommets.

\origpage{324}

Nous pourrions également envisager le polyèdre $P'$, défini plus haut, et construire deux tableaux relatifs l'un aux faces de $P'$, l'autre à ses arêtes.

Comparons le tableau (1), relatif aux faces de $P$, avec le tableau (1~\emph{bis}) relatif aux arêtes de $P'$.

Il résulte de ce qui précède, que ces deux tableaux peuvent se déduire l'un de l'autre en remplaçant les lignes par les colonnes et inversement.

Cela posé, envisageons le tableau (1), relatif aux faces de $P$, et examinons comment on peut déduire de ce tableau le nombre de \emph{Betti} $P_{2}$ du polyèdre $P$.

\emph{Comment, d'abord, pourra-t-on en déduire les congruences entre les faces et les arêtes}?

Considérons une colonne quelconque de la 1\textsuperscript{ère} sorte; par exemple la $i^{\text{ème}}$ colonne. Multiplions les éléments de cette colonne et de la $k^{\text{ème}}$ ligne de la 1\textsuperscript{ère} sorte par $a^{2}_{k}$ et ajoutons; puis égalons à la somme obtenue, en multipliant les élements de cette même colonne et de la $j^{\text{ème}}$ ligne de la 2\textsuperscript{de} sorte par $a^{1}_{j}$; nous obtiendrons la congruence
\[ a^{2}_{i} \equiv \sum \varepsilon^{2}_{i,j} a^{1}_{j}, \]
ce qui est bien une des congruences (3) du \S~II. Toutes les autres congruences n'en sont que des combinaisons.

\emph{Comment pourra-t-on maintenant trouver les homologies entre les faces}?

Pour cela, envisageons par exemple la $i^{\text{ème}}$ colonne de la 2\textsuperscript{de} sorte; multiplions les éléments de la $k^{\text{ème}}$ ligne et de cette colonne par $a^{2}_{k}$, ajoutons et égalons à zéro; nous trouverons~:
\[ \sum \varepsilon^{3}_{i,k} a^{2}_{k} \backsim 0, \]
ce qui est bien une des homologies (5) du \S~II, dont toutes les autres ne sont que des combinaisons.

\emph{Qu'adviendra-t-il, maintenant, si l'on applique à notre tableau} (1) \emph{une transformation algébrique quelconque}?

Avant la transformation, chaque colonne de la 1\textsuperscript{ère} sorte correspond
\origpage{325}
à une face, chaque colonne de la 2\textsuperscript{de} sorte à une case, chaque ligne de la 1\textsuperscript{ère} sorte à une face, chaque ligne de la 2\textsuperscript{de} sorte à une arête.

On obtient, comme nous l'avons vu, autant de congruences et d'homologies que de colonnes, en multipliant les éléments de chaque ligne de la 1\textsuperscript{ère} sorte par la face correspondante, ceux de chaque ligne de la 2\textsuperscript{de} sorte par l'arête correspondante, et ajoutant.

Supposons maintenant qu'on fasse la deuxième opération, c'est-à-dire qu'on transforme en ajoutant la ligne de la 1\textsuperscript{ère} sorte, qui correspond à $a^{2}_{i}$, à celle de la 1\textsuperscript{ère} sorte, qui correspond à $a^{2}_{k}$. Nous convenons de dire qu'à la nouvelle $k^{\text{ème}}$ ligne (celle à laquelle on a ajouté la $i^{\text{ème}}$ ligne) correspond toujours la variété $a^{2}_{k}$; mais qu'à la nouvelle $i^{\text{ème}}$ ligne (qui d'ailleurs n'a pas changé) correspond la variété $a^{2}_{i} - a^{2}_{k}$.

Si l'on fait la cinquième opération sur la $k^{\text{ème}}$ ligne de la 1\textsuperscript{ère} sorte, en en multipliant les éléments par une constante $m$, nous conviendrons de dire qu'à la nouvelle $k^{\text{ème}}$ ligne correspond la variété $\dfrac{1}{m} a^{2}_{k}$ (notation qui n'a qu'une valeur symbolique, à moins que $\dfrac{1}{m}$ ne soit entier).

Quant à la quatrième opération, ce n'est qu'une combinaison de plusieurs opérations analogues à la deuxième.

Nous avons ainsi défini la variété qui correspond à chacune des lignes de la 1\textsuperscript{ère} sorte du tableau, après qu'on a appliqué à ces lignes une combinaison quelconque des 2\textsuperscript{ème}, 4\textsuperscript{ème} et 5\textsuperscript{ème} opérations.

Nous définirions de même les variétés qui correspondent aux différentes lignes de la 2\textsuperscript{de} sorte, après qu'on aurait appliqué à ces lignes une combinaison des 2\textsuperscript{ème}, 4\textsuperscript{ème} et 5\textsuperscript{ème} opérations.

Grâce à ces conventions, il suffira encore, pour obtenir les congruences et les homologies, d'ajouter et d'égaler à zéro, après avoir multiplié les éléments de chaque ligne par la variété correspondante, et avoir changé le signe des produits ainsi obtenus, en ce qui concerne les lignes de la 2\textsuperscript{de} sorte.

Maintenant, si l'on applique aux colonnes du tableau les 1\textsuperscript{ème}, 3\textsuperscript{ème} et 5\textsuperscript{ème} opérations, on ne fera que combiner les congruences entre elles, et les homologies entre elles; ou multiplier une congruence et une homologie par un facteur constant.

D'où le résultat suivant~:

\origpage{326}

\emph{Pour déduire les congruences du tableau transformé, voici ce qu'il faut faire}~: multiplier chaque ligne de la 1\textsuperscript{ère} sorte pour la variété qui lui correspond en vertu de la convention que nous venons de faire, et ajouter; faire de même pour les lignes de la 2\textsuperscript{de} sorte; égaler les deux résultats ainsi obtenus; on aura ainsi autant de congruences que de colonnes de la 1\textsuperscript{ère} sorte; toutes les autres congruences possibles ne seront que des combinaisons.

\emph{Pour déduire de même les homologies, il faut}~: multiplier chaque ligne de la 1\textsuperscript{ère} sorte pour la variété correspondante, ajouter et égaler à zéro; on aura ainsi autant d'homologies que de colonnes de la 2\textsuperscript{de} sorte; toutes les autres homologies possibles n'en seront que des combinaisons.

Il importe de remarquer que les congruences et homologies ainsi obtenues, pourront n'avoir qu'une valeur symbolique, parce que les coefficients pourront être fractionnaires.

Et, en effet, d'une part les éléments du tableau transformé peuvent ne plus être entiers; d'autre part, la variété qui correspond à une ligne peut, comme je l'ai dit plus haut, n'avoir elle-même qu'une valeur symbolique.

Mais comme les coefficients, entiers ou non, sont toujours commensurables, il suffira de multiplier notre congruence ou notre homologie par un entier convenable, pour en déduire une congruence ou une homologie à coefficients entiers, qui aura un sens pour elle-même.

\emph{Supposons, maintenant, qu'on ait réduit le tableau, comme je l'ai dit plus haut}.

Combien y aura-t-il d'homologies distinctes?

Parmi nos $N_{3}$ colonnes de la 2\textsuperscript{de} sorte, il y en aura $N_{3} - N'_{2}$ dont tous les éléments seront nuls, et $N_{2}$\ednote{Printed $N_{2}$; the context requires $N'_{2}$, an apparent misprint.} dont un élément sera égal à 1 et tous les autres nuls. Les $N_{3} - N'_{2}$ premières ne nous donneront aucune homologie; chacune des $N'_{2}$ autres nous en donnera une et ces $N'_{2}$ homologies seront évidemment toutes distinctes.

Il y a donc $N'_{2}$ homologies distinctes.

Combien y a-t-il de congruences distinctes entre les faces et les arêtes?

Il y en a évidemment $N_{2}$, correspondant aux $N_{2}$ colonnes de la
\origpage{327}
1\textsuperscript{ère} sorte, et ces congruences sont distinctes, parce que le déterminant formé avec les lignes et les colonnes de la 1\textsuperscript{ère} sorte, n'est pas nul.

Considérons maintenant dans notre tableau réduit les $N_{1}$ lignes de la 2\textsuperscript{de} sorte; parmi elles il y en aura $N_{1} - N''_{2}$ dont tous les éléments sont nuls, et $N''_{2}$, dont un élément est égal à 1 et tous les autres nuls. Parmi nos $N_{2}$ congruences il y en aura donc $N''_{2}$ qui contiendront une arête, et $N_{2} - N'_{2}$\ednote{Printed $N'_{2}$ with a single accent; the congruences containing no edge number $N_{2} - N''_{2}$, as the next sentence states, so a misprint or a lost second accent.} qui ne contiendront aucune arête. Il y a donc $N_{2} - N''_{2}$ congruences \emph{entre les faces seulement}, et ces congruences sont toutes distinctes.

Il y aura donc \emph{entre les faces seulement} $N_{2} - N'_{2} - N''_{2}$ congruences, qui resteront distinctes, si l'on ne regarde plus comme distinctes celles que l'on peut déduire les unes des autres par le moyen des homologies.

\emph{Le nombre de Betti relatif aux faces de $P$ est donc}~:
\[ N_{2} - N'_{2} - N''_{2} + 1. \]

Cherchons maintenant le nombre de \emph{Betti} relatif aux arêtes de $P$.

On le trouvera évidemment en opérant comme nous venons de faire sur le tableau (1~\emph{bis}), relatif aux arêtes de $P'$.

Mais on passe d'un tableau à l'autre, en remplaçant les lignes par les colonnes, et réciproquement. Les nombres qui joueront, par rapport à (1~\emph{bis}), le même rôle que
\[ N_{2}, \qquad N'_{2}, \qquad N''_{2} \]
jouent par rapport à (1), seront donc respectivement~:
\[ N_{2}, \qquad N''_{2}, \qquad N'_{2}. \]

Donc le nombre de \emph{Betti} relatif aux arêtes de $P'$ est encore~:
\[ N_{2} - N'_{2} - N''_{2} + 1. \]

Ainsi les nombres de \emph{Betti} relatifs, l'un aux faces de $P$, l'autre aux arêtes de $P'$, sont égaux.

\origpage{328}

Or nous avons vu plus haut que les nombres de \emph{Betti} relatifs aux arêtes de $P$ et à celles de $P'$ sont égaux, de même que les nombres de \emph{Betti} relatifs aux faces de $P$ et à celles de $P'$.

\emph{Donc le nombre de Betti relatif aux faces de $P$ est égal au nombre de Betti relatif aux arêtes de $P$}.

Notre théorème fondamental est donc démontré en ce qui concerne le polyèdre $P$, c'est-à-dire, en ce qui concerne les polyèdres de l'espace à quatre dimensions.

La démonstration pourrait, sans aucun doute, s'étendre à un polyèdre quelconque.

\subsection*{\S~IX. Remarques diverses.}

Le théorème fondamental est ainsi établi par une démonstration, qui diffère essentiellement de celle de la page 43 de l'\emph{Analysis situs}.

Mais cela ne saurait pas nous suffire. Il faut nous efforcer de retrouver les propositions intermédiaires, et, en particulier, celle-ci~:

La condition nécessaire et suffisante pour que l'on puisse trouver une variété $V$ telle que $\displaystyle\sum N(V, V_{i}) \neq 0$, c'est que l'on n'ait pas l'homologie $\displaystyle\sum V_{i} \backsim 0$.

Considérons deux variétés, la première $V_{1}$, à une dimension, composée d'arêtes de $P'$, la seconde $V_{2}$, à deux dimensions, composée de faces de $P$, de telle sorte qu'on aura~:
\[ V_{1} = \sum \alpha_{i} b^{1}_{i}, \qquad V_{2} = \sum \alpha'_{i} a^{2}_{i}, \]
l'arête $b^{1}_{i}$ étant celle qui correspond à la face $a^{2}_{i}$, d'après les conventions du \S~7.

L'arête $b^{1}_{i}$ coupe la face $a^{2}_{i}$, et n'en coupe aucune autre, de telle sorte que si nous reprenons la notation de l'\emph{Analysis situs}, pag.\ 38, nous aurons~:
\[ N(V_{1}, V_{2}) = \sum \alpha_{i} \alpha'_{i}. \]

Nous supposerons dans ce qui va suivre, que les variétés $V_{1}$ et
\origpage{329}
$V_{2}$ sont fermées, ce qui s'exprime par les congruences~:
\[ \sum \alpha_{i} b^{1}_{i} \equiv 0, \qquad \sum \alpha'_{i} a^{2}_{i} \equiv 0. \tag{1} \]

Vérifions d'abord que l'on aura
\[ \sum \alpha_{i} \alpha'_{i} = 0, \]
pourvu que l'on ait l'une des deux homologies${}^{(*)}$~:
\[ \sum \alpha_{i} b^{1}_{i} \backsim 0, \qquad \sum \alpha'_{i} a^{2}_{i} \backsim 0. \tag{2} \]

Si en effet nous avons, par exemple, la seconde homologie (2), c'est qu'on aura~:
\[ \alpha'_{i} = \sum_{j=1}^{j=N_{3}} \zeta_{j} \varepsilon^{3}_{i,j}, \]\ednote{Printed $\varepsilon^{3}_{i,j}$; the argument that follows and (4) use $\varepsilon^{3}_{j,i}$, an apparent misprint.}
$\zeta_{j}$ étant un coefficient ne dépendant que de $j$.

D'un autre côté, la première des congruences (1) peut se déduire de l'une des suivantes~:
\[ b^{1}_{i} \equiv \sum b^{0}_{j} \varepsilon^{3}_{j,i}, \tag{3} \]
d'où
\[ \sum \alpha_{i} b^{1}_{i} \equiv \sum \sum \alpha_{i} b^{0}_{j} \varepsilon^{3}_{j,i}. \]

En égalant à zéro le coefficient de $b^{0}_{j}$, il vient successivement~:
\[ \sum \alpha_{i} \varepsilon^{3}_{j,i} = 0, \qquad \sum \sum \alpha_{i} \zeta_{j} \varepsilon^{3}_{j,i} = 0, \qquad \sum \alpha_{i} \alpha'_{i} = 0, \qquad \text{C.~Q.~F.~D.} \]

On raisonnerait de même si on avait la première homologie (2).

Je dis maintenant que si la seconde homologie (2) n'a pas lieu, on peut choisir les $\alpha_{i}$ de telle façon que $V_{1}$ reste fermée et cependant que.\ednote{Printed ``que.'' with a stray point.} $\displaystyle\sum \alpha_{i} \alpha'_{i}$ ne soit pas nul.

\textbf{(*)} Cfr.\ \emph{Analysis situs}, pag.\ 42.

En effet, dire que la seconde homologie (2) n'a pas lieu, c'est
\origpage{330}
dire qu'on ne peut pas trouver des nombres $\zeta_{j}$ tels que l'on ait~:
\[ \alpha'_{i} = \sum \zeta_{j} \varepsilon^{3}_{j,i}. \tag{4} \]

Dire que $V_{1}$ reste fermée, c'est dire que les $\alpha_{i}$ sont assujettis aux conditions~:
\[ \sum \alpha_{i} \varepsilon^{3}_{j,i} = 0. \tag{5} \]

Or il est clair que si les $\alpha'_{i}$ ne satisfont pas à des égalités de la forme (4), l'équation linéaire $\displaystyle\sum \alpha_{i} \alpha'_{i} = 0$ sera distincte des équations (5); on pourra donc toujours trouver des nombres $\alpha_{i}$, qui satisfassent aux équations (5) sans satisfaire à $\displaystyle\sum \alpha_{i} \alpha'_{i} = 0$.

Remarquons d'ailleurs que nous n'avons pas restreint la généralité en supposant que nos variétés $V_{1}$ et $V_{2}$ étaient des combinaisons des $b^{1}_{i}$ et des $a^{2}_{i}$, quelle que soit la variété $V$, dont la subdivision forme les polyèdres $P$ et $P'$. Quelles que soient les variétés $V_{1}$ et $V_{2}$, nous pouvons toujours subdiviser $V$, de manière à former deux polyèdres réciproques $P$ et $P'$ tels que $V_{1}$ soit une combinaison des arêtes du second, et $V_{2}$ une combinaison des faces du premier.

Il faudrait voir comment le tableau (1) du \S~8 et les tableaux analogues, peuvent nous permettre de déterminer les nombres de \emph{Betti}, tels que \emph{Betti} les définit lui-même, et non plus les nombres de \emph{Betti} définis de la seconde manière, c'est-à-dire ceux que nous avons considérés jusqu'à présent.

Considérons, par exemple, un tableau analogue au tableau (1), mais relatif aux arêtes du polyèdre $P$ et à leurs relations avec les faces et les sommets. Considérons, en particulier, les colonnes de la 2\textsuperscript{de} sorte et les signes\ednote{Printed ``signes'' for ``lignes''; an apparent misprint.} de la 1\textsuperscript{ère} sorte, où figurent les nombres $\varepsilon^{2}_{i,j}$. Soit $T$ le tableau partiel ainsi obtenu. A l'aide de ce tableau on pourra former les congruences
\[ a^{2}_{i} \equiv \sum \varepsilon^{2}_{i,j} a^{1}_{j}, \]
d'où l'on déduit les homologies
\[ \sum \varepsilon^{2}_{i,j} a^{1}_{j} \backsim 0. \tag{6} \]

Alors pour reconnaître si plusieurs lignes fermées, formées de
\origpage{331}
combinaisons des arêtes $a^{1}_{j}$ sont distinctes \emph{au sens de la 1\textsuperscript{ère} définition}, c'est-à-dire, \emph{au sens de Betti}, il faut savoir si elles sont liées par une homologie obtenue en combinant les homologies (6) par addition, soustraction ou multiplication, mais \emph{sans division}.

Supposons qu'on ait appliqué à notre tableau une série de ces transformations, que j'ai appelées arithmétiques au \S~8.

Soit $\zeta^{2}_{i,j}$ le nombre qui, dans le tableau transformé, figure dans la $j^{\text{ème}}$ ligne de la 1\textsuperscript{ère} sorte et la $i^{\text{ème}}$ colonne de la 2\textsuperscript{de} sorte. Soit $c_{j}$ la variété qui correspond à la $j^{\text{ème}}$ ligne de la 1\textsuperscript{ère} sorte de notre tableau transformé en vertu des conventions du \S~8. D'après ce que nous avons vu dans ce \S~8, cette variété n'est qu'une combinaison des arêtes $a^{1}_{j}$. Nous aurons alors les homologies
\[ \sum \zeta^{2}_{i,j} c_{j} \backsim 0. \tag{6 \textit{bis}} \]

Ces homologies ne sont que des combinaisons des homologies (6), que l'on peut obtenir \emph{sans division}, et réciproquement on peut tirer des homologies (6) des homologies (6~\emph{bis}) \emph{sans division}, c'est là une conséquence du caractère arithmétique des transformations.

On peut donc, quand on veut s'assurer si deux lignes fermées sont distinctes au sens de \emph{Betti}, se servir des homologies (6~\emph{bis}) au lieu des homologies (6).

Nous pouvons supposer qu'on s'est servi des transformations arithmétiques pour réduire le tableau, comme je l'ai expliqué au \S~8, et par conséquent que $\xi^{2}_{i,j}$\ednote{Printed $\xi$; the number defined above is $\zeta^{2}_{i,j}$, an apparent misprint.} est nul~: 1° si $i > j$; 2° si $j > N_{2}$.

Le tableau réduit aux colonnes de la 2\textsuperscript{de} sorte et aux lignes de la 1\textsuperscript{ère} sorte prendra, par exemple, la forme suivante~:
\[ \begin{matrix} a & 0 & 0 & 0 & 0 \\[1ex] e & b & 0 & 0 & 0 \\[1ex] f & g & c & 0 & 0 \\[1ex] h & k & l & d & 0 \\[1ex] 0 & 0 & 0 & 0 & 0 \\[1ex] 0 & 0 & 0 & 0 & 0. \end{matrix} \]

\origpage{332}

J'ai supposé 6 lignes et 5 colonnes; j'ai supposé que l'un des nombres $\zeta^{2}_{i,i}$ est égal à zéro, de façon qu'une des colonnes du tableau transformé soit entièrement composée de zéros. J'ajoute que si $d$ était égal à 1, les nombres $h$, $k$, $l$, qui figurent à la même ligne, seraient nuls.

Cela posé, si $d$ n'est pas égal à 1, les deux définitions des nombres de \emph{Betti} ne coïncident pas, parce que l'on a l'homologie $dc_{4} \backsim 0$, d'où l'on ne pourrait déduire l'homologie $c_{4} \backsim 0$ que par division. Si $d = 1$, on a $h = k = l = 0$, et si $c$ n'est pas égal à 1, on aura l'homologie $cc_{3} \backsim 0$, et les deux définitions ne concorderont pas; et ainsi de suite.

En résumé, pour que les deux définitions concordent, il faut et il suffit que le produit $abcd$ soit égal à 1.

Pour interpréter ce résultat, revenons au tableau non transformé. Le produit $abcd$ sera le plus grand commun diviseur de tous les déterminants obtenus en supprimant $N_{2} - N_{3}$ lignes dans le tableau $T$, pourvu que ces déterminants ne soient pas tous nuls (auquel cas il n'y aurait pas dans le tableau transformé de colonne exclusivement composée de zéros). Si les déterminants sont tous nuls, on en formera d'autres en supprimant dans le tableau $T$ une colonne et $N_{2} - N_{3} + 1$ lignes; le produit $abcd$ sera le plus grand commun diviseur de tous ces déterminants, s'ils ne sont pas tous nuls; et ainsi de suite.

Nous arrivons ainsi à la règle suivante.

Soit $\Delta_{p}$ le plus grand commun diviseur des déterminants obtenus en partant du tableau $T$ et y supprimant $p$ lignes et $N_{2} - N_{3} + p$ colonnes.\ednote{The rule deletes $p$ lignes and $N_{2} - N_{3} + p$ colonnes, while the preceding paragraph deletes $N_{2} - N_{3}$ lignes, then one colonne and $N_{2} - N_{3} + 1$ lignes; the two do not agree as printed.} La condition nécessaire et suffisante pour que les deux définitions des nombres de \emph{Betti} coïncident, c'est que le premier des $\Delta_{p}$ qui ne s'annule pas, soit égal à 1 (le plus grand commun diviseur de plusieurs nombres égaux à zéro étant zéro par définition).

Supposons que la variété $V_{1} = \displaystyle\sum \alpha_{i} b^{1}_{i}$, considérée au début de ce paragraphe, ne peut pas être frontière d'une variété à deux dimensions, mais satisfait à l'homologie $V_{1} \backsim 0$. En d'autres termes, l'homologie $V_{1} \backsim 0$ peut se déduire des homologies (6) par division, mais non pas sans division. Dans ce cas, on aura néanmoins~:
\[ N(V_{1}, V_{2}) = \sum \alpha_{i} \alpha'_{i} = 0. \]

\origpage{333}

\subsection*{\S~X. Démonstration arithmétique de l'un des théorèmes du \S~VII.}

Voici une manière de former les homologies qui pourra être utile à connaître.

Soit $b^{0}_{i}$ un sommet du polyèdre $P'$, situé à l'intérieur d'une case $a^{3}_{i}$ du polyèdre $P$. Soit, d'autre part, $a^{0}_{k}$ un sommet de $P$, appartenant à la case $a^{3}_{i}$. Joignons $b^{0}_{i}$ à $a^{0}_{k}$ par une ligne que j'appellerai simplement $b^{0}_{i} a^{0}_{k}$.

Soit maintenant $b^{1}_{i}$ une arête de $P'$, dont les deux extrémités sont $b^{0}_{j}$ et $b^{0}_{h}$, de telle sorte que l'une des congruences (3) (cfr.\ \S~2) relatives à $P'$ soit~:
\[ b^{1}_{i} \equiv b^{0}_{j} - b^{0}_{h}. \]

Soit, d'autre part, $a^{2}_{i}$ la face de $P$ qui correspond à l'arête $b^{1}_{i}$ de $P'$, et $a^{0}_{k}$ un des sommets de $a^{2}_{i}$; nous aurons l'homologie~:
\[ b^{1}_{i} \backsim a^{0}_{k} b^{0}_{j} - a^{0}_{k} b^{0}_{h}. \tag{1} \]

Soit $a^{1}_{i}$ une arête de $P$, dont les deux extrémités sont $a^{0}_{j}$ et $a^{0}_{h}$, de sorte que l'une des congruences (3) relatives à $P$ soit~:
\[ a^{1}_{i} \equiv a^{0}_{j} - a^{0}_{h}. \]

Soit $a^{3}_{k}$ l'une des cases de $P$ auxquelles appartient $a^{1}_{i}$, et $b^{0}_{k}$ le sommet correspondant de $P'$; on aura l'homologie~:
\[ a^{1}_{i} \backsim b^{0}_{k} a^{0}_{j} - b^{0}_{k} a^{0}_{h}. \tag{2} \]

Je dis maintenant que toutes les homologies entre les $a^{1}_{i}$ peuvent se déduire des homologies (2).

En effet, soit $a^{2}_{i}$ une face quelconque de $P$, et soit~:
\[ a^{2}_{i} \equiv \sum \varepsilon^{2}_{i,j} a^{1}_{j} \]
\origpage{334}
la congruence de la forme (3) qui lui correspond; on en déduit l'homologie~:
\[ \sum \varepsilon^{2}_{i,j} a^{1}_{j} \backsim 0, \tag{3} \]
et nous avons vu au \S~6 que toutes les homologies entre les arêtes de $P$ sont des combinaisons de celles qu'on obtient de la sorte.

Soit alors $a^{1}_{j}$ l'une des arêtes de $P$ qui figurent dans l'homologie (3), et soit~:
\[ a^{1}_{i} \equiv a^{0}_{h} - a^{0}_{l}. \]
\ednote{Printed $a^{1}_{i}$ on the left; the edge just introduced, and the homology (2 bis) below, call it $a^{1}_{j}$, so $a^{1}_{j}$ appears to be meant.}

Soit d'ailleurs $a^{3}_{k}$ l'une des cases dont fait partie $a^{2}_{i}$. Nous aurons l'homologie~:
\[ a^{1}_{j} \backsim b^{0}_{k} a^{0}_{h} - b^{0}_{k} a^{0}_{l}. \tag{2 \textit{bis}} \]

Si nous additionnons les homologies (2~\emph{bis}) qui sont de la forme (2), après les avoir multipliées par $\varepsilon^{2}_{i,j}$, tous les termes du second membre disparaîtront en vertu des relations (5) du \S~2; on retrouverait donc l'homologie (3). C.~Q.~F.~D.

On démontrerait de même que toutes les homologies entre les $b^{1}_{i}$ peuvent se déduire des homologies (1).

Nous avons vu plus haut, au \S~7, que si l'on a une congruence~:
\[ \sum a^{1}_{i} \equiv 0, \]
on peut trouver une autre congruence entre les arêtes de $P'$~:
\[ \sum b^{1}_{j} \equiv 0, \]
et de telle façon qu'on ait l'homologie~:
\[ \sum a^{1}_{i} \backsim \sum b^{1}_{j}. \tag{4} \]

Je dis maintenant que cette homologie (4) peut être déduite des homologies (1) et (2).

Découpons, en effet, le premier membre de notre congruence $\displaystyle\sum a^{1}_{i} \equiv 0$ en un certain nombre de groupes, de telle façon que les
\origpage{335}
arêtes d'un même groupe appartiennent à une même case $a^{3}_{k}$. Soit $\displaystyle\sum a^{1}_{j}$ l'un de ces groupes; nous aurons la congruence~:
\[ \sum a^{1}_{j} \equiv a^{0}_{m} - a^{0}_{p}, \tag{5} \]
$a^{0}_{m}$ et $a^{0}_{p}$ étant les deux extrémités de la ligne formée par l'ensemble des arêtes de ce groupe. Je suppose que toutes ces arêtes appartiennent à la case $a^{3}_{k}$. Soit
\[ a^{1}_{j} \equiv a^{0}_{h} - a^{0}_{l} \]
l'une de ces arêtes; nous aurons l'homologie~:
\[ a^{1}_{j} \equiv b^{0}_{k} a^{0}_{h} - b^{0}_{k} a^{0}_{l}, \tag{2 \textit{ter}} \]
\ednote{The print sets the congruence sign $\equiv$ here, though the text calls this an homology and it is of the form (2); the homology sign $\backsim$ appears to be meant.}
et en ajoutant toutes ces homologies, on trouverait~:
\[ \sum a^{1}_{j} \backsim b^{0}_{k} a^{0}_{m} - b^{0}_{k} a^{0}_{p}. \tag{6} \]

Ajoutons d'une part toutes les homologies (6), d'autre part toutes les congruences (5), qui correspondent aux différents groupes. L'addition des congruences (5) doit nous donner la congruence $\displaystyle\sum a^{1}_{i} \equiv 0$; il s'en suit que si un sommet $a^{0}_{m}$ figure dans une des congruences (5) avec le signe $+$, il devra figurer dans une autre avec le signe $-$. L'addition des homologies (6) nous donnera donc~:
\[ \sum a^{1}_{j} \backsim \sum (b^{0}_{k} a^{0}_{m} - b^{0}_{q} a^{0}_{m}). \tag{7} \]

En écrivant cette relation je suppose que $a^{0}_{m}$ figure dans deux des congruences (5), une fois avec le signe $+$ dans la congruence qui correspond à la case $a^{3}_{k}$, et une fois avec le signe $-$ dans la congruence qui correspond à la case $a^{3}_{q}$.

Observons maintenant que $b^{0}_{k}$ et $b^{0}_{q}$ sont deux sommets de $P'$, et que ces deux sommets appartiennent l'un et l'autre à la case $b^{3}_{m}$. On peut alors trouver une ligne formée d'arêtes de $P'$, appartenant à cette case $b^{3}_{m}$, et allant de $b^{0}_{k}$ à $b^{0}_{q}$. Soit $\displaystyle\sum b^{1}_{s}$ cette ligne; on aura
\[ \sum b^{1}_{s} \equiv b^{0}_{k} - b^{0}_{q}. \tag{5 \textit{bis}} \]

\origpage{336}

De même que de la congruence (5) des homologies (2~\emph{ter}), qui sont de la forme (2), nous avons déduit l'homologie (6); de même de la congruence (5~\emph{bis}) et d'homologies de la forme (1), nous pourrons déduire l'homologie~:
\[ \sum b^{1}_{s} \backsim a^{0}_{m} b^{0}_{k} - a^{0}_{m} b^{0}_{q}. \tag{6 \textit{bis}} \]

A chaque terme du second membre de (7) correspond une homologie (6~\emph{bis}). En les additionnant, on trouvera~:
\[ \sum \sum b^{1}_{s} \backsim \sum (a^{0}_{m} b^{0}_{k} - a^{0}_{m} b^{0}_{q}), \tag{7 \textit{bis}} \]
d'où~:
\[ \sum a^{1}_{j} + \sum \sum b^{1}_{s} \backsim 0, \tag{8} \]
homologie de la forme (4), qui se déduit, comme on le voit, des homologies (1) et (2). C.~Q.~F.~D.

On peut se demander pourquoi j'ai jugé nécessaire de revenir sur un théorème déjà démontré au \S~7. On le comprendra si on se rend compte de la nature géométrique, pour ainsi dire, de la démonstration du \S~7. La présente démonstration a, au contraire, un caractère arithmétique; elle n'invoque que des propriétés des schémas définis au \S~2, et des tableaux construits au \S~8; et elle conserverait sa valeur alors même qu'à ces schémas et à ces tableaux ne correspondrait aucun polyèdre.

Qu'avons-nous supposé en effet? C'est que si $\alpha^{p}_{0}, \alpha^{p}_{1}, \alpha^{p}_{2}$ sont les nombres des sommets, des arêtes et des faces appartenant à une même case, et si $\beta^{p}_{0}, \beta^{p}_{1}, \beta^{p}_{2}$ sont les nombres des cases, des faces et des arêtes auxquelles appartient un même sommet, on a~:
\[ \alpha^{p}_{0} - \alpha^{p}_{1} + \alpha^{p}_{2} = \beta^{p}_{0} - \beta^{p}_{1} - \beta^{p}_{2} = 2; \]
\ednote{The print sets a minus before $\beta^{p}_{2}$; the Euler relation for the cells around a vertex, parallel to the left-hand member, requires $\beta^{p}_{0} - \beta^{p}_{1} + \beta^{p}_{2}$. The dash shows no trace of a lost vertical stroke, so this is reproduced as an apparent misprint.}
et en outre que deux sommets quelconques $a^{0}_{i}$ et $a^{0}_{k}$ sont liés par l'homologie~:
\[ a^{0}_{i} \backsim a^{0}_{k}. \tag{9} \]

Or on peut reconnaître si un sommet appartient à une face, par
\origpage{337}
exemple, en appliquant au tableau du \S~8 des règles purement arithmétiques, et on peut de la même manière reconnaître si une homologie telle que (9) a lieu.

\subsection*{\S~XI. Possibilité de la subdivision.}

Tout ce qui précède suppose qu'une variété quelconque peut être subdivisée en variétés simplement connexes, de manière à former un polyèdre $P$, à $p$ dimensions, pour lequel les variétés $a^{p}_{i}, a^{p-1}_{i}, \ldots, a^{2}_{i}, a^{1}_{i}, a^{0}_{i}$ sont toutes simplement connexes. Par exemple, toute variété à trois dimensions pourra être subdivisée en cases simplement connexes, séparées les unes des autres par des faces simplement connexes.

C'est cela qu'il nous reste à démontrer, et c'est cette démonstration que je vais donner. Je précise davantage~: je vais montrer que toute variété à $p$ dimensions peut être subdivisée de façon à former un polyèdre $P$, dont toutes les variétés $a^{p}_{i}, a^{p-1}_{i}, \ldots, a^{2}_{i}, a^{1}_{i}, a^{0}_{i}$ sont des tétraèdres généralisés.

Je supposerai que le théorème a été démontré pour une variété à $p - 1$ dimensions, et je me propose de l'étendre à une variété à $p$ dimensions.

Nous présenterons la définition de notre variété sous la forme suivante, qui comprendra les deux définitions données dans l'\emph{Analysis situs}.

Nous aurons les équations et les inégalités
\[ \begin{aligned} &x_{i} = \theta_{i}(y_{1}, y_{2}, \ldots, y_{q}), \qquad (i = 1, 2, \ldots, n)\\ &f_{k}(y_{1}, y_{2}, \ldots, y_{q}) = 0, \qquad (k = 1, 2, \ldots, q - p)\\ &\varphi_{h}(y_{1}, y_{2}, \ldots, y_{q}) > 0. \end{aligned} \tag{1} \]

Ces équations et ces inégalités définiront une variété $v$ qui sera limitée et, en général, non fermée; on aura différents systèmes analogues d'équations et d'inégalités, définissant autant de variétés partielles que j'appellerai $v_{1}, v_{2}, \ldots, v_{m}$.

\origpage{338}

Deux de ces variétés seront dites contigües, si elles ont une partie commune, et je puis supposer que l'on peut passer d'un point quelconque de l'une de ces variétés à un point quelconque d'une autre quelconque d'entre elles, sans sortir de l'ensemble de ces variétés. Cet ensemble constituera la variété que j'appellerai $V$, et qu'il s'agissait de définir.

Je supposerai que cette variété $V$ est bilatère.

C'est évidemment là la façon la plus générale possible de définir une variété.

Considérons la variété partielle $v_{1}$, définie par les équations (1).

D'après le théorème des fonctions implicites, on pourra satisfaire aux équations
\[ f_{k} = 0, \]
en faisant
\[ y_{j} = \psi_{j}(z_{1}, z_{2}, \ldots, z_{p}), \]
les $\psi$ étant des fonctions holomorphes des $z$; mais les séries $\psi$ pourront ne pas converger pour tous les points de la variété $v_{1}$.

Les conditions de convergence seront certaines inégalités~:
\[ \eta_{k}(z_{1}, z_{2}, \ldots, z_{p}) > 0. \]

Quand on remplacera les $y$ en fonctions des $z$, les relations~:
\[ x_{i} = \theta_{i}, \qquad \varphi_{h} > 0 \]
deviendront~:
\[ \begin{aligned} &x_{i} = \theta'_{i}(z_{1}, z_{2}, \ldots, z_{p})\\ &\varphi'_{h}(z_{1}, z_{2}, \ldots, z_{p}) > 0. \end{aligned} \]

Alors l'ensemble des relations
\[ x_{i} = \theta'_{i}, \qquad \varphi'_{h} > 0, \qquad \eta_{k} > 0 \tag{1 \textit{bis}} \]
définira une certaine variété $v'_{1}$, de telle façon que l'ensemble des variétés analogues à $v'_{1}$ reproduira la variété $v_{1}$.

\origpage{339}

Nous sommes ainsi ramenés à la seconde définition de l'\emph{Analysis situs}.

Cela posé, soit $v''_{1}$ une autre variété $v'_{1}$ satisfaisant aux conditions suivantes~: elle sera tout entière contenue dans $v'_{1}$; elle comprendra tous les points de $v'_{1}$ qui ne lui sont pas communs avec une des variétés contigües; par conséquent la frontière complète de $v''_{1}$ sera tout entière dans la partie commune à $v'_{1}$ et aux variétés contigües.

A chacune des variétés
\[ v'_{1}, \qquad v'_{2}, \qquad \ldots, \]
dont l'ensemble constitue $V$, correspondra ainsi une variété
\[ v''_{1}, \qquad v''_{2}, \qquad \ldots, \]
satisfaisant aux conditions que je viens d'énoncer; et il est clair qu'on peut s'arranger de telle façon que tout point de $V$ appartienne à l'une des variétés $v''$, et à une seule, à moins qu'il ne soit sur la frontière de l'une des variétés $v''$, auquel cas il devra appartenir, en outre, à la frontière au moins d'une autre variété $v''$.

La variété $V$, ainsi subdivisée en variétés $v''$, constitue un polyèdre $P$, au sens donné à ce mot au \S~2. Mais ce polyèdre ne convient pas encore à la question, car nous ne pouvons savoir si les variétés $v''$ sont des tétraèdres généralisés, ou même sont simplement connexes.

Considérons la variété $v''_{1}$, et soit~:
\[ z_{1} = 0, \qquad z_{2} = 0, \qquad \ldots, \qquad z_{p} = 0 \]
un point \emph{intérieur à cette variété}; considérons la variété à une dimension~:
\[ z_{1} = \alpha_{1} t, \qquad z_{2} = \alpha_{2} t, \qquad \ldots, \qquad z_{p} = \alpha_{p} t, \]
où les $\alpha$ sont des constantes, et où nous ferons varier $t$ depuis 0 jusqu'à $+\infty$. C'est ce que j'appellerai un \emph{rayon vecteur}.

Chaque rayon vecteur rencontrera la frontière complète de $v''_{1}$ en un nombre impair de points; en effet, quand on suivra ce rayon, en
\origpage{340}
faisant varier $t$ de 0 à $+\infty$, on sortira de la variété $v''_{1}$; on pourra y rentrer ensuite et en sortir plusieurs fois, mais on finira toujours par en sortir pour n'y plus rentrer.

Il pourra se faire qu'un rayon vecteur rencontre la frontière de $v''_{1}$ en deux points confondus. Les rayons vecteurs qui satisfont à cette condition s'appelleront les \emph{rayons remarquables}.

L'ensemble des rayons remarquables formera une ou plusieurs variétés à $p - 1$ dimensions, que j'appellerai les \emph{cônes remarquables}.

Les intersections des cônes remarquables avec la frontière de $v''_{1}$ formeront une ou plusieurs variétés à $p - 2$ dimensions, que j'appellerai $U$, et ces variétés $U$ partageront la frontière de $v''_{1}$ en régions que j'appellerai $R$.

Une région $R$ ne peut être rencontrée par un rayon vecteur en plus d'un point; mais d'après ce que nous venons de voir, il peut se présenter deux cas~: quand on suit ce rayon vecteur, en faisant croître $t$ de 0 à $\infty$, on peut, au moment où on rencontre $R$, sortir de $v''_{1}$ ou y rentrer. Si le premier cas, par exemple, se présente pour un des vecteurs qui rencontrent $R$, il se présentera pour tous les vecteurs qui rencontrent $R$.

D'où la distinction des régions $R$ en régions de la 1\textsuperscript{ère} sorte, que les rayons vecteurs rencontrent en sortant de $v''_{1}$, et en régions de la 2\textsuperscript{de} sorte, que les rayons vecteurs rencontrent en rentrant dans $v''_{1}$.

Les régions $R$ étant des variétés à $p - 1$ dimensions pourront, d'après l'hypothèse faite au début, être subdivisées en tétraèdres généralisés.

Supposons, pour fixer les idées, qu'un rayon vecteur rencontre trois fois la frontière de $v''_{1}$, qu'il rencontre successivement les régions $R_{1}, R_{2}, R_{3}$; $R_{1}$ et $R_{3}$ seront de la 1\textsuperscript{ère} sorte, $R_{2}$ sera de la 2\textsuperscript{de} sorte.

Subdivisons $R_{1}$ et $R_{3}$ en tétraèdres généralisés à $p - 1$ dimensions.

Si $T_{1}$ est une des subdivisions de $R_{1}$, menons tous les rayons vecteurs qui passent par les différents points de $T_{1}$, et conservons la partie de ces rayons vecteurs qui est comprise entre le point $z_{i} = 0$ et le rayon\ednote{Sic in the print. $R_{1}$ is a region of the frontier of $v''_{1}$, not a ray; ``la région'' is presumably meant.} $R_{1}$ (partie qui est intérieure à $v''_{1}$); l'ensemble de ces vecteurs formera un tétraèdre généralisé à $p$ dimensions, ayant pour sommet le point $z_{i} = 0$, et pour base le tétraèdre généralisé à $p - 1$ dimensions $T_{1}$.

\origpage{341}

Soit maintenant $T_{3}$ une des subdivisions de $R_{3}$; menons encore tous les rayons vecteurs qui passent par les différents points de $T_{3}$ et conservons la partie de ces rayons vecteurs qui est comprise entre $R_{2}$ e\ednote{Printed ``e'', apparently a misprint for ``et''.} $R_{3}$ (partie qui est intérieure à $v''_{1}$). Cet ensemble forme une variété à $p - 1$ dimensions que l'on pourrait appeler un \emph{tronc de tétraèdre généralisé}, dont les deux bases sont $T_{3}$ et un tétraèdre généralisé à $p - 1$ dimensions, que j'appellerai $T_{2}$ et qui fera partie de $R_{2}$. C'est, en d'autres termes, la différence de deux tétraèdres généralisés, ayant pour sommet commun le point $z_{i} = 0$ et pour bases l'un $T_{3}$, l'autre $T_{2}$.

Ce tronc de tétraèdre généralisé pourra à son tour être partagé en $p$ tétraèdres généralisés, de même que dans le théorème classique, le tronc de pyramide triangulaire se partage en trois pyramides triangulaires.

Finalement $v''_{1}$ sera partagé en tétraèdres généralisés.

Une difficulté subsiste cependant; on peut subdiviser comme $v''_{1}$ les autres variétés analogues $v''_{2}, v''_{3}, \ldots$; considérons la subdivision de $v''_{1}$ en tétraèdres généralisés $T_{1}$ et celle de $v''_{2}$ en tétraèdres généralisés $T_{2}$. La frontière commune de $v''_{1}$ et $v''_{2}$ se trouvera subdivisée d'une part en tétraèdres généralisés à $p - 1$ dimensions $\tau_{1}$, qui seront les faces des $T_{1}$, et d'autres part\ednote{Printed ``d'autres part'', a misprint for ``d'autre part''.} en tétraèdres généralisés à $p - 1$ dimensions $\tau_{2}$, qui seront les faces des $T_{2}$; \emph{mais il n'est pas évident que ces deux subdivisions coïncident}.

Considérons alors la partie commune à l'un des $\tau_{1}$ et à l'un des $\tau_{2}$; je pourrai, d'après l'hypothèse faite au début, la subdiviser en tétraèdres généralisés à $p - 1$ dimensions $\sigma$. Ainsi chacun des tétraèdres $\tau_{1}$ et chacun des tétraèdres $\tau_{2}$ sera subdivisé en tétraèdres $\sigma$.

Soit maintenant $\tau'_{1}$ une des variétés à $q$ dimensions \emph{appartenant} à $\tau_{1}$ (j'emploie ici le mot appartenir dans le même sens que quand je dis que les faces, les arêtes et les sommets d'un tétraèdre ordinaire appartiennent à ce tétraèdre, ou quand je disais au \S~2 que les variétés $a^{q}_{i}$ appartenaient au polyèdre $P$). Soit de même $\tau'_{2}$ une des variétés à $q$ dimensions appartenant à $\tau_{2}$. Ces deux variétés $\tau'_{1}$ et $\tau'_{2}$ seront des tétraèdres généralisés, puisque d'après la définition du tétraèdre généralisé, toute variété qui \emph{appartient} à un tétraèdre généralisé, est elle-même un tétraèdre généralisé. Alors $\tau'_{1}$ et $\tau'_{2}$ se trouveront subdivisés en tétraèdres
\origpage{342}
généralisés à $q$ dimensions $\sigma$\ednote{Printed $\sigma$; the text below refers to these tetrahedra as ``ceux que j'ai appelés $\sigma'$'', so $\sigma'$ appears to be meant.}, qui \emph{appartiendront} aux tétraèdres à $p - 1$ dimensions $\sigma$.

Cela à la rigueur pourrait nous suffire; nos variétés $v''_{1}$, etc., seraient partagées en tétraèdres généralisés à $p$ dimensions $T^{p}$, leurs frontières en tétraèdres à $p - 1$ dimensions $T^{p-1}$, etc.; seulement ces tétraèdres $T^{p-1}$ ne seraient pas ceux qui appartiennent aux tétraèdres $T^{p}$, cela en seraient seulement des subdivisions.

Mais nous pouvons aller plus loin.

Considérons l'un des tétraèdres à $p$ dimensions $T^{p}$ dans lequel $v''_{1}$ est subdivisé. Je rappelle qu'on les a obtenus en subdivisant les troncs de tétraèdres généralisés, dont il a été question plus haut. Par conséquent $T^{p}$ a tous ses sommets sur la frontière de $v''_{1}$ (il y aurait exception pour les tétraèdres dont un sommet est au point $z_{i} = 0$, mais pour ceux-là il n'y a pas de difficulté).

Supposons, par exemple, que les points communs à $T^{p}$ et à la région que j'ai appelée plus haut $R_{3}$ forment un tétraèdre généralisé à $q$ dimensions $T^{q}$ \emph{appartenant} à $T^{p}$, et que les points communs à $T^{p}$ et à la région $R_{2}$ forment un tétraèdre à $p - q - 1$ dimensions $T^{p-q-1}$, appartenant à $T^{p}$.

Les tétraèdres $T^{q}$ et $T^{p-q-1}$ sont analogues aux tétraèdres $\tau'_{1}$ traités plus haut; ils peuvent donc être subdivisés en tétraèdres analogues à ceux que j'ai appelés $\sigma'$; soient $S^{q}_{1}, S^{q}_{2} \ldots$, les tétraèdres analogues à $\sigma'$, qui sont des subdivisions de $T^{q}$; soient $S^{p-q-1}_{1}, S^{p-q-1}_{2}, \ldots$ les tétraèdres analogues à $\sigma'$ qui sont des subdivisions de $T^{p-q-1}$. Je dis qu'on peut subdiviser $T^{p}_{k}$ en tétraèdres à $p$ dimensions de telle façon que les variétés $S^{q}_{1}, S^{q}_{2}, \ldots, S^{p-q-1}_{1}, S^{p-q-1}_{2}, \ldots$ appartiennent à $T^{p}$.

Pour le démontrer je suppose d'abord que $T^{p}$ soit un \emph{tétraèdre rectiligne} (cfr.\ \S~2 in fine). On sait qu'un tétraèdre rectiligne est entièrement défini quand on connait ses $p + 1$ sommets. Alors $T^{p}$ est le tétraèdre rectiligne qui a pour sommets ceux de $T^{q}$ et de $T^{p-q-1}$.

Supposons que $T^{q}$ se décompose en $g$ tétraèdres partiels~:
\[ S^{q}_{1}, \qquad S^{q}_{2}, \qquad \ldots, \qquad S^{q}_{g}, \]
et $T^{p-q-1}$ en $h$ tétraèdres partiels~:
\[ S^{p-q-1}_{1}, \qquad S^{p-q-1}_{2}, \qquad \ldots, \qquad S^{p-q-1}_{h}. \]

\origpage{343}

On vérifiera alors que $T^{p}$ se décompose en $gh$ tétraèdres partiels qui sont ceux dont les sommets sont ceux de
\[ S^{q}_{i} \text{ et } S^{p-q-1}_{k} \qquad (i = 1, 2, \ldots, g;\ k = 1, 2, \ldots, h). \]

Si le tétraèdre $T^{p}$ n'est pas rectiligne, le résultat subsiste puisque un tétraèdre quelconque est homéomorphe à un tétraèdre rectiligne.

Ainsi notre variété est décomposée en tétraèdres à $p$ dimensions de façon à former un polyèdre tel que toute variété appartenant à ce polyèdre, appartient à l'un de ces tétraèdres.

On est ainsi débarrassé des derniers doutes qui pouvaient subsister au sujet de la possibilité de subdiviser une variété $V$ de façon à former un polyèdre $P$, pour lequel tous les $a^{q}_{i}$ soient simplement connexes.

Paris, mars 1899.

H.~Poincaré.

\end{document}
